\documentclass[11pt,oneside]{article}

\usepackage[T1]{fontenc}
\usepackage[utf8]{inputenc}
\usepackage[letterpaper,margin=1in]{geometry}
\usepackage{amsmath,amssymb}
\usepackage{siunitx}
\usepackage{booktabs}
\usepackage{longtable}
\usepackage{tabularx}
\usepackage{array}
\usepackage{enumitem}
\usepackage{graphicx}
\usepackage{microtype}
\usepackage{xcolor}
\usepackage{listings}
\usepackage{tikz}
\usetikzlibrary{arrows.meta,backgrounds,calc,fit,positioning,shapes.geometric}
\usepackage{natbib}
\usepackage{hyperref}
\usepackage{xurl}
\usepackage[nameinlink,noabbrev]{cleveref}

\definecolor{modsimblue}{HTML}{235789}
\definecolor{modsimcyan}{HTML}{2A9D8F}
\definecolor{modsimorange}{HTML}{E76F51}
\definecolor{modsimgold}{HTML}{E9C46A}
\definecolor{modsimgray}{HTML}{F3F5F7}
\definecolor{codegray}{HTML}{44515C}

\hypersetup{
  colorlinks=true,
  linkcolor=modsimblue,
  citecolor=modsimcyan,
  urlcolor=modsimblue,
  pdftitle={ModSim: Architecture, Physics Runtime, and SMORES-EP Demonstrations},
  pdfauthor={ModSim Project Contributors}
}

\lstdefinestyle{modsim}{
  basicstyle=\ttfamily\small,
  backgroundcolor=\color{modsimgray},
  frame=single,
  rulecolor=\color{black!15},
  breaklines=true,
  columns=fullflexible,
  keepspaces=true,
  showstringspaces=false,
  keywordstyle=\color{modsimblue}\bfseries,
  commentstyle=\color{codegray}\itshape
}
\lstset{style=modsim}

\newcommand{\modsim}{\textsc{ModSim}}
\newcommand{\robotpack}{Robot Pack}
\newcommand{\mujoco}{MuJoCo}
\newcommand{\smoresep}{SMORES-EP}
\newcommand{\code}[1]{\texttt{#1}}
\newcommand{\pathcode}[1]{\texttt{\detokenize{#1}}}
\newcommand{\yes}{\textcolor{modsimcyan}{Implemented}}
\newcommand{\partial}{\textcolor{modsimgold!70!black}{Partial}}
\newcommand{\future}{\textcolor{modsimorange}{Future}}
\newcommand{\sourcefact}{\textsuperscript{P}}
\newcommand{\cadfact}{\textsuperscript{G}}
\newcommand{\simchoice}{\textsuperscript{S}}
\newcommand{\measuredfact}{\textsuperscript{M}}

\newcolumntype{Y}{>{\raggedright\arraybackslash}X}
\newcolumntype{L}[1]{>{\raggedright\arraybackslash}p{#1}}
\setlist{nosep,leftmargin=*}
\setcounter{secnumdepth}{3}
\setcounter{tocdepth}{2}

\title{\textbf{ModSim: A Backend-Agnostic Semantic and Physics Framework\\
for Modular and Multi-Robot Simulation}\\[0.7em]
\large Architecture, Robot-Pack Authoring, Runtime Docking, Model Views,\\
and Wheel-Driven SMORES-EP Reconfiguration}
\author{ModSim Project Contributors}
\date{Technical Report --- Implementation Snapshot 0.1.0\\
Repository commit \texttt{4a5ca00} --- \today}

\begin{document}

\maketitle

\begin{abstract}
Modular and multi-robot research requires more than rigid-body dynamics. A useful
simulation framework must also represent platform semantics: which physical bodies
constitute a reusable module, where docking interfaces are located, which interfaces
are compatible, what geometric and dynamic conditions constitute an acceptable dock,
how a physical constraint becomes a logical connection, and how changing connections
alter the robot's topology. Existing physics engines are powerful at integration,
contact, and rendering, but they do not provide these modular-robot concepts as a
portable, backend-independent contract. Conversely, platform-specific research code
often embeds topology, controller, file-layout, and simulator assumptions into a
single application, making validation, reuse, and cross-backend development difficult.

This report presents \modsim, a Python-first framework that separates canonical
modular-robot semantics from simulator-specific mechanics. The framework introduces a
strict, self-contained \robotpack{} format; a desktop Robot Pack Studio; a canonical
event-driven runtime; deterministic docking acceptance and two-phase physical/logical
commit; an adapter boundary with mock and \mujoco{} implementations; immutable
model-view generation; and a Runtime Inspector that couples a native three-dimensional
viewer with a live topology graph, event log, status, and metrics. The implementation
is evaluated with generic synthetic fixtures and a committed \smoresep{} package
containing detailed visual geometry, a physics-specific MJCF model, four EP-face
connectors, effort-controlled joints, wheel/ground contact, and runtime welds.

Three progressively more demanding demonstrations are analyzed. A two-module physical
cycle drives one module through differential-wheel effort, docks its TOP/PAN face to a
rear BOTTOM face, holds an idealized weld, releases it after an 80-ms delay, and reverses
away. A seven-module kinematic Driver-to-Snake demonstration exercises topology and
events. A seven-module physics realization executes the same four published topology
replacements using gravity, tire contact, measured connector frames, authored
collision-clearance routes, and coherent wheel commands for moving three-module
assemblies. The reference trace performs four $6\rightarrow5\rightarrow6$ connection
transitions and reaches the final seven-module chain with ten successful docks, four
successful undocks, and no semantic failures. The recorded software verification
baseline contains 537 passing tests with 86\% branch-aware core coverage.

The report emphasizes epistemic boundaries. Published hardware facts, geometry derived
from the Fusion export, provisional simulation parameters, and measured simulation
outcomes are identified separately. The current work demonstrates an operational
semantic and physics vertical slice; it does not yet claim calibrated motor dynamics,
magnetic capture physics, autonomous reconfiguration planning, hardware trajectory
reproduction, or backend equivalence beyond the tested semantic contract.
\end{abstract}

\noindent\textbf{Keywords:} modular robotics, self-reconfiguration, robot simulation,
MuJoCo, SMORES-EP, docking, topology graph, digital twin, robot description, runtime
visualization.

\vspace{1em}
\noindent\textbf{Evidence notation.} In tables, \sourcefact{} denotes a published
platform fact, \cadfact{} a value derived from committed geometry, \simchoice{} a
provisional simulation or controller choice, and \measuredfact{} a result observed in
the reference simulation or enforced by a regression test.

\tableofcontents
\clearpage

\section{Introduction}
\label{sec:introduction}

Modular robots promise a form of embodied adaptability that a conventional
fixed-morphology robot cannot provide. A collection of repeated modules can change
topology, redistribute actuation, form task-specific kinematic chains, recover from
failures, or separate into independently mobile agents. The configuration of the
robot is therefore not merely a parameter of its controller; it is a time-varying
part of the system state. This observation has shaped the field of modular
self-reconfigurable robotics for decades \citep{yim2007modular}.

Simulation infrastructure has not always followed the same abstraction. General
robotics engines accept bodies, joints, actuators, collision geometry, and constraints.
Those primitives are necessary, but they do not answer modular-robot questions such
as:

\begin{itemize}
  \item Which links and joints constitute one reusable module type?
  \item Which physical surfaces are connectors, and which pairs may mate?
  \item What translational, angular, roll, and speed errors are acceptable at capture?
  \item When is an engine constraint considered a committed logical connection?
  \item Which modules form an assembly after a dock or release?
  \item How should a changing morphology be exposed to planners, algorithms, logs, and
        graphical tools?
\end{itemize}

In a simulator-specific prototype, these concepts are often distributed across model
files, Python scripts, controller constants, scene-construction code, and assumptions
inside a viewer. That approach can produce an effective demonstration, but it creates
three long-term problems. First, semantic knowledge is not portable: replacing one
engine with another requires reconstructing the meaning of connectors and
configurations. Second, physical and logical state can diverge: an application may add
an edge to its topology graph before the engine has actually created a constraint, or
may remove a graph edge while a stale constraint remains active. Third, the tooling
surface becomes fragile: authoring, validation, simulation, algorithms, and
visualization each invent a slightly different representation of the same robot.

\modsim{} addresses this gap with a semantics-first architecture. It is not a new
rigid-body solver. Instead, it defines a portable \robotpack{}, a canonical runtime,
and narrow interfaces to physics engines and user interfaces. The central design rule
is:

\begin{quote}
\emph{ModSim owns modular-robot meaning and canonical history; a backend owns physics;
every visualization or algorithmic view is derived from the canonical state.}
\end{quote}

This separation makes the framework useful even before multiple high-fidelity backends
exist. A dependency-free mock adapter can exercise connector semantics and event
ordering in continuous integration. A \mujoco{} adapter can then supply articulated
dynamics, contact, effort actuation, runtime welds, and native visualization without
moving connector rules into engine-specific code. Future adapters---for example an
Isaac Sim implementation---can be checked against the same semantic contract.

\subsection{Research questions}

The current implementation is organized around four questions.

\begin{enumerate}
  \item \textbf{Semantic portability.} Can one platform description define modules,
        connectors, docking policy, metadata, and generated model views without
        depending on one simulator's scene format?
  \item \textbf{Runtime consistency.} Can physical constraints, logical connections,
        assembly membership, events, and graph views remain coherent during docking
        and undocking?
  \item \textbf{Progressive physical fidelity.} Can a generic semantic runtime begin
        with kinematic tests and then support real gravity, contact, articulated joint
        control, and multi-module locomotion without replacing its state model?
  \item \textbf{Inspectability.} Can authors and runtime users validate the same data,
        observe the same event history, and view physical and topological state from a
        single authoritative simulation?
\end{enumerate}

\subsection{Contributions}

This technical report describes four principal contributions.

\begin{enumerate}
  \item A strict, self-contained \robotpack{} format that separates mechanical assets
        from modular semantics, backend mappings, custom metadata, capabilities, and
        model-view recipes. The associated loader, validator, canonical writer, URDF
        importer, and Studio application provide an end-to-end authoring workflow.
  \item A backend-independent runtime that resolves connector frames, evaluates
        compatibility and geometric acceptance, applies policy guards, performs a
        two-phase physical/logical commit, evolves assemblies, accepts atomic joint
        command batches, and records deterministic semantic events and metrics.
  \item Narrow physics and visualization boundaries: a reference mock backend, an
        optional \mujoco{} backend, immutable model-view data-transfer objects, a
        revision-aware view factory, and a dual-window Runtime Inspector driven by one
        authoritative session.
  \item A \smoresep{} case study containing committed visual and collision assets,
        four docking faces, a physics-specific MJCF model, wheel and joint effort
        control, tire contact, a physical two-module docking cycle, and a seven-module
        wheel-driven realization of a published Driver-to-Snake topology plan.
\end{enumerate}

\subsection{Scope and non-claims}

The implementation is pre-alpha software (version 0.1.0), and its strongest evidence
is software and simulation evidence. The report does \emph{not} claim a new contact
solver, autonomous reconfiguration planner, calibrated \smoresep{} digital twin,
magnetic-field model, or reproduction of a published Driver-to-Snake hardware
trajectory. The seven-module routes, friction parameters, effort limits, skid model,
and control gains are explicit simulation choices. The published 2019 work supplies
the topology and connector substitutions; \modsim{} supplies one physically plausible,
deterministic execution of that plan.

\subsection{Organization}

\Cref{sec:background} reviews relevant modular-robot and simulation concepts.
\Cref{sec:requirements} turns the observed gap into requirements and authority rules.
\Cref{sec:architecture} presents the overall architecture. The following four sections
describe Robot Packs and Studio, runtime docking, backend physics, and model views with
runtime visualization. \Cref{sec:smores} develops the \smoresep{} case study and its
controllers. \Cref{sec:experiments,sec:analysis} report experiments and analyze the
evidence. \Cref{sec:future} states limitations and a staged roadmap. The appendices
provide schema summaries, event taxonomies, scenario commands, and a reproducibility
checklist.


\section{Background, Related Work, and the Addressed Gap}
\label{sec:background}

\subsection{Modular self-reconfigurable robots}

Modular self-reconfigurable robot systems combine repeated hardware units with
mechanisms for changing connectivity. Their potential benefits include functional
reuse, task-dependent morphology, redundancy, repair, and scalable manufacturing;
their enduring challenges include mechanical connection, configuration planning,
control, communication, and the combinatorial growth of possible morphologies
\citep{yim2007modular}. A configuration is naturally represented by a graph whose
vertices denote modules and whose edges denote connections, but a useful physical
description must also preserve connection faces, orientation, joint state, spatial
frames, and actuation.

SMORES and \smoresep{} are mobile-style modular robots whose left and right faces also
serve as wheels. The original SMORES design demonstrated locomotion and
reconfiguration \citep{davey2012smores}. \smoresep{} adds electro-permanent-magnet
connectors. Each module has four active rotational degrees of freedom---LEFT, RIGHT,
PAN, and TILT---and four mating faces---LEFT, RIGHT, TOP, and BOTTOM
\citep{liu2020parallel,liu2023smoresep,modlabSmoresEp}. LEFT and RIGHT provide
differential-drive locomotion; LEFT, RIGHT, and PAN rotate continuously, while TILT is
bounded.

The EP-Face was designed as a planar, genderless, switchable connector. Published
characterization reports high holding force, rapid state switching, and tolerance to
position error \citep{tosun2016epface}. These properties motivate a simulation model
with explicit connector frames, capture conditions, release timing, and post-capture
constraints. They also illustrate why a bare URDF joint tree is insufficient: a
connector's compatibility, allowable roll orientations, capture region, and release
policy are not standard URDF concepts.

\subsection{Planning topology versus executing motion}

Liu, Whitzer, and Yim formulate reconfiguration as a sequence of connector changes and
demonstrate a distributed planning approach \citep{liu2019distributed}. Their
Driver-to-Snake example gives an initial and target configuration and four connector
replacement actions. This is valuable topology-level ground truth, but it does not
fully specify motor torques, contact coefficients, obstacle-clearance routes, or a
time-parameterized trajectory for the entire Driver-to-Snake sequence. Later SMORES
self-assembly work develops navigation, pose adjustment, and straight-approach control
and demonstrates hardware behaviors \citep{liu2020parallel,liu2023smoresep}.

The distinction between \emph{what connections should change} and \emph{how bodies
should move} is central to \modsim{}. A \code{ReconfigurationPlan} describes connector
replacements. A physical scenario supplies platform mechanics, control parameters, and
routes. The semantic runtime remains responsible for validating and committing the
connection. This separation allows future planners to replace authored routes without
changing the docking or event model.

\subsection{Physics engines and robot-description formats}

Modern simulators already solve difficult numerical problems: articulated dynamics,
contact, equality constraints, actuators, sensors, and visualization. \mujoco{} is
designed for efficient model-based control and represents multijoint dynamics in
generalized coordinates \citep{todorov2012mujoco}. Its native MJCF language can express
rich dynamic models, and its compiler can also import URDF
\citep{mujocoModelingDocs,urdfFormat}. Isaac Sim and related USD-based workflows offer
another path to high-fidelity scene composition and rendering. Gazebo/SDF, Webots, and
CoppeliaSim provide still other combinations of physics, scene, and control tooling.

These engines intentionally operate below the domain model proposed here. A weld
constraint can hold two bodies together, but it does not by itself establish:

\begin{itemize}
  \item that two surfaces are semantically compatible;
  \item that their approach axes and roll orientations satisfy a connector contract;
  \item that a runtime command or auto-latch policy allows connection;
  \item that the logical module graph should gain a particular stable edge; or
  \item that downstream algorithms should receive a canonical dock event.
\end{itemize}

URDF has a complementary limitation. It is widely supported as a mechanical
description of links, joints, inertial data, and geometry, but it does not standardize
modular-robot docking semantics. MJCF can encode engine-specific contact and actuator
details, but using MJCF alone would make the portable semantic package engine-specific.
\modsim{} therefore treats URDF/MJCF as mechanical assets referenced by a separate
semantic package.

\subsection{Why an engine-independent semantic layer is needed}

The addressed gap lies between a robot asset and a modular-robot application. The asset
describes what can move and collide. The application needs stable identities,
compatibility rules, topology, lifecycle events, validation, and generated views.
Embedding those concepts directly in a backend creates several failure modes:

\begin{enumerate}
  \item \textbf{Duplicated truth.} A graph, a scene graph, and a controller each carry a
        different list of active docks.
  \item \textbf{Optimistic commit.} The application records a dock even though the
        engine refuses or cannot allocate a physical constraint.
  \item \textbf{Frame drift.} A connector pose authored in one tool is recomposed
        differently by another tool, so visual alignment and physical acceptance
        disagree.
  \item \textbf{Non-portable scenes.} Connector semantics are hidden in simulator
        scripts and cannot be validated before that simulator is installed.
  \item \textbf{Uninspectable failures.} A boolean docking result omits which criterion
        failed and by how much.
  \item \textbf{Unsafe extensibility.} A data package contains arbitrary executable
        plugin paths rather than stable builder or backend identifiers.
\end{enumerate}

\subsection{Positioning relative to existing simulators}

\Cref{tab:positioning} summarizes the intended division of responsibility. The table
does not claim that general simulators cannot be extended to provide these features;
rather, it identifies the contract \modsim{} adds above their native physical model.

\begin{table}[htbp]
\centering
\caption{Conceptual positioning of \modsim{} relative to a physics engine.}
\label{tab:positioning}
\begin{tabularx}{\textwidth}{L{0.25\textwidth}YY}
\toprule
Concern & Physics engine or robot asset & \modsim{} semantic layer \\
\midrule
Rigid-body state & Integrates generalized coordinates, contact, and constraints &
Ingests observations into stable module/link/connector identities \\
Mechanical description & URDF, MJCF, or a future USD asset &
References assets and adds portable module semantics \\
Dock validity & Can test distance or create a raw constraint &
Defines compatibility, acceptance, policy, lifecycle, and failure detail \\
Dynamic connection & Activates/deactivates an equality or joint &
Uses two-phase commit before mutating logical topology \\
Configuration & Scene/body graph & Canonical module multigraph and assemblies \\
History & Engine-dependent logging & Ordered semantic event log and metrics \\
Visualization & Native 3-D scene & Derived, immutable topology and event views \\
Authoring & Engine-specific editor or XML & Schema-aware Robot Pack Studio \\
\bottomrule
\end{tabularx}
\end{table}

The resulting architecture is intentionally compositional. \modsim{} benefits from a
mature solver instead of reimplementing one; the solver benefits from a structured
domain layer that can be tested independently. A future C++ execution core, browser
client, or second simulator can reuse the same semantic contracts.


\section{Requirements, System Model, and Design Principles}
\label{sec:requirements}

\subsection{Terminology and entities}

The following distinctions are used throughout the report.

\begin{description}[style=nextline]
  \item[Module type and module instance]
  A module type is reusable semantic and mechanical data in a \robotpack{}. A module
  instance is one placed occurrence in a runtime scene.

  \item[Connector type and connector instance]
  A connector type defines compatibility, gender, capture tolerance, orientation,
  physical constraint intent, and policy. A connector instance attaches that meaning to
  a link and local pose on a module type.

  \item[Physical constraint and logical connection]
  A physical constraint is a backend object such as a \mujoco{} weld. A logical
  connection is canonical \modsim{} state created only after the physical constraint
  has been accepted.

  \item[Assembly]
  An assembly is a connected component of the committed module graph. A disconnected
  module is an assembly of size one.

  \item[Canonical state and model view]
  \code{WorldState} is mutable runtime authority. A model view is an immutable,
  JSON-safe projection generated from that authority for algorithms or interfaces.

  \item[Reconfiguration plan and physical route]
  A plan names connector removals and additions. A route describes how a mobile module
  or assembly attempts to realize one plan action through physics.
\end{description}

\subsection{Functional requirements}

\begin{longtable}{L{0.10\textwidth}L{0.31\textwidth}L{0.47\textwidth}}
\caption{Primary requirements and their implemented mechanisms.}
\label{tab:requirements}\\
\toprule
ID & Requirement & Implemented mechanism \\
\midrule
\endfirsthead
\toprule
ID & Requirement & Implemented mechanism \\
\midrule
\endhead
R1 & Portable platform package & Strict split-YAML Robot Pack with relative assets,
stable identifiers, metadata, mappings, and recipes. \\
R2 & Early and profile-aware validation & Authoring and simulation-readiness profiles
with structured issues and cross-reference checks. \\
R3 & Backend independence & Lazy adapter registry, capability negotiation, and no
mandatory simulator dependency in core. \\
R4 & Deterministic docking & Sorted candidates, structured acceptance, guards,
explicit commands, stable connection identifiers, and deterministic events. \\
R5 & Physical/logical consistency & Two-phase commit: backend creates the constraint
before \code{DockCommitted} mutates canonical state. \\
R6 & Dynamic topology & Event-driven connection changes, assembly recomputation,
revision counters, and generated topology multigraphs. \\
R7 & Actuated dynamics & Backend-neutral atomic joint commands with declaration,
capability, finiteness, and limit validation; MuJoCo effort implementation. \\
R8 & Inspectable execution & Immutable Runtime Inspector frames, contiguous event
deltas, graph/status/metrics UI, and native physics viewer from one session. \\
R9 & Safe authoring & Local assets, path containment, strict schemas, transactional
updates, canonical writing, and non-overwriting export. \\
R10 & Extensibility & Registries and typed protocols for backends and model-view
builders without executable code embedded in Robot Packs. \\
\bottomrule
\end{longtable}

\subsection{Non-functional requirements}

\paragraph{Determinism.}
Stable IDs, sorted iteration, canonical YAML, explicit event sequence numbers, and
deterministic candidate ordering make failures reproducible. Physical trajectories may
still depend on solver and floating-point behavior, but semantic arbitration should not
depend on hash or insertion order.

\paragraph{Dependency isolation.}
The base package depends only on Pydantic, Rich, ruamel.yaml, and Typer. \mujoco{} and
Studio reside behind independent optional extras. Importing \code{modsim} must not
import Qt, VTK, or \mujoco{}.

\paragraph{Fail-closed behavior.}
Unsupported physical connection types, absent backend capabilities, exhausted weld
pools, unresolved frames, invalid command modes, and bad cross-references should
produce actionable errors rather than a simulated success.

\paragraph{Bounded memory.}
The event log is intentional history, while model-view caches retain only the latest
result per source and recipe. UI publication is throttled and transfers immutable
snapshots rather than sharing backend objects.

\paragraph{Auditability.}
Every committed dock and undock has a canonical event. Runtime logs are local,
per-launch sidecars. The GUI does not silently persist edits: users explicitly Save or
Export a canonical pack.

\subsection{Authority rules}

The most important architectural decisions can be expressed as authority rules.

\begin{enumerate}
  \item A mechanical asset owns detailed geometry and kinematics. Robot Pack YAML
        should not duplicate an entire link tree.
  \item A Robot Pack owns portable modular semantics. A backend model should not be the
        only place a docking face is named or typed.
  \item \code{WorldState} owns live logical state. A graph widget or backend scene must
        not become an independent connection database.
  \item A backend owns integration and the existence of physical constraints. Core
        code never writes directly into \mujoco{} arrays.
  \item Events own committed semantic history. Candidates and transient acceptance
        state may be ephemeral; topology mutations may not.
  \item Model views own presentation-oriented projections only. They can be rebuilt
        from the pack and world and are never edited to control canonical state.
\end{enumerate}

\subsection{Current non-goals}

The first phase intentionally excludes a new physics solver, a CAD system, a general
autonomous planner, a calibrated magnetic connector model, arbitrary plugin discovery,
and a web application. The APIs are designed so these can be added without replacing
the semantic core. This is particularly important for a future high-performance C++
runtime: Python coordinates authoring and research workflows today, while immutable
data contracts and narrow backend protocols avoid making Python object identity part of
the long-term wire model.


\section{Overall Architecture}
\label{sec:architecture}

\subsection{Layered organization}

\modsim{} is distributed as \code{modsim-robotics} and currently supports Python 3.11
and newer. The source tree is divided into three packages:

\begin{itemize}
  \item \pathcode{src/modsim}: backend-neutral schemas, loading, validation, canonical
        runtime, docking, model views, scenarios, CLI, and mock backend;
  \item \pathcode{src/modsim_backend_mujoco}: optional scene composition, adapter,
        viewer, runtime host process, and weld/contact mechanics; and
  \item \pathcode{src/modsim_studio}: optional Qt/PyVista authoring and Qt/PyQtGraph
        runtime visualization.
\end{itemize}

\Cref{fig:layers} shows the authority and dependency direction. Upward clients depend
on immutable or typed core interfaces; the core does not import optional UI or physics
packages.

\begin{figure}[htbp]
\centering
\begin{tikzpicture}[
  node distance=5mm and 8mm,
  box/.style={draw,rounded corners,minimum height=9mm,align=center,font=\small,
              fill=modsimgray,text width=0.28\textwidth},
  core/.style={box,fill=modsimblue!12,draw=modsimblue},
  opt/.style={box,fill=modsimcyan!12,draw=modsimcyan!80!black},
  data/.style={box,fill=modsimgold!20,draw=modsimgold!70!black},
  arrow/.style={-{Latex[length=2mm]},thick,draw=black!65}
]
  \node[data] (assets) {Mechanical assets\\URDF / MJCF / meshes};
  \node[data,right=of assets] (pack) {Robot Pack\\semantics / mappings / recipes};
  \node[core,below=of $(assets)!0.5!(pack)$,text width=0.60\textwidth] (load)
    {Loader, validator, writer, importer, scene specification};
  \node[core,below=of load,text width=0.60\textwidth] (runtime)
    {Canonical runtime: WorldState, docking, assemblies, events, metrics, commands};
  \node[opt,below left=of runtime] (backends) {Backend adapters\\mock / MuJoCo / future};
  \node[core,below=of runtime] (views) {ModelViewFactory\\immutable derived views};
  \node[opt,below left=of views] (studio) {Robot Pack Studio\\authoring};
  \node[opt,below right=of views] (inspector) {Runtime Inspector\\graph / events / metrics};
  \draw[arrow] (assets) -- (load);
  \draw[arrow] (pack) -- (load);
  \draw[arrow] (load) -- (runtime);
  \draw[arrow] (runtime) -- (backends);
  \draw[arrow] (backends) -- node[below,sloped,font=\scriptsize]{snapshots / outcomes} (runtime);
  \draw[arrow] (runtime) -- (views);
  \draw[arrow] (load) -| (studio);
  \draw[arrow] (views) -- (inspector);
\end{tikzpicture}
\caption{Layered \modsim{} architecture. Arrow direction denotes the primary flow of
data or requests, not Python import direction. Optional physics and UI packages remain
outside the dependency-light core.}
\label{fig:layers}
\end{figure}

\subsection{End-to-end data flow}

An ordinary workflow moves through six stages.

\begin{enumerate}
  \item The user imports a concrete URDF or opens an existing pack. The loader resolves
        split YAML documents and local assets under a pack root.
  \item Validation checks schema and semantic cross-references under an authoring or
        simulation-readiness profile.
  \item A \code{SceneSpec} instantiates module types with stable instance IDs and initial
        poses.
  \item A backend loads the selected mechanical asset for each instance and returns a
        handle registry. The runtime immediately ingests a snapshot.
  \item Each step advances physics, ingests observed state, processes overload and
        docking, appends events, and updates revision counters.
  \item A model-view factory builds an immutable topology projection, and the Runtime
        Inspector publishes it with event deltas, metrics, and scenario status.
\end{enumerate}

The crucial property is that no stage invents a second canonical topology. The backend
reports constraints and frames; the world stores logical connections; the graph is
derived from those connections.

\subsection{Stable identity}

Typed string identifiers are used for pack entities and runtime instances. A connector
instance ID combines a module instance and local connector, for example
\code{module\_3/pan}. A joint instance similarly combines module and semantic joint.
A connection ID is canonical with respect to endpoint ordering, so requesting the same
undirected pair in either order yields the same identity. Assembly IDs derive from
membership, using the smallest module ID in the component, rather than from a process
counter. These rules support deterministic logging, graph keys, replay, and
cross-process transport.

\subsection{Capability-oriented extension}

A backend advertises exactly what it implements. The base contract includes loading,
stepping, snapshots, physical connection creation/removal, shutdown, and capabilities.
Optional protocols add module pose/twist control or joint commands. Core code checks
capabilities before use. For example, a backend without runtime constraints produces a
failed dock rather than a logical-only edge; a backend without effort control rejects a
physical wheel scenario during preflight, before the initial world is mutated.

Model-view extensibility follows the same pattern. A pack recipe names a stable builder
ID, not a Python import path. An application registers builder implementations in a
\code{ModelViewFactory}. Missing builders are valid portable data but unavailable in
that installation, producing an explicit construction error.

\subsection{Package and distribution boundary}

The wheel contains the three Python packages. The source distribution additionally
contains tests, examples, and documentation. This creates one known portability
boundary: the named \smoresep{} demonstration configurations currently live in
\pathcode{examples/scenarios} and are intended for a source checkout. The installed
library exposes the generic scenario machinery, but a wheel without the example tree
cannot execute those named source-checkout demos. A future scenario resource or
explicit \code{--scenario} mechanism should remove this mismatch.

\subsection{Implementation scale}

At the current snapshot, the three source packages contain 72 Python modules and about
20,455 lines of Python. The test directory contains 40 test modules and about 11,073
lines. The committed \smoresep{} pack occupies approximately 16 MiB, dominated by five
binary STL visual meshes. These numbers characterize engineering scope, not runtime
performance; \Cref{sec:analysis} treats performance separately.


\section{Robot Packs, Validation, Persistence, and Studio}
\label{sec:robotpacks}

\subsection{Robot Pack as the portable unit}

A \robotpack{} is a directory whose root document is
\pathcode{robot_pack.yaml}. It contains split semantic specifications, backend mapping
documents, and declared local assets. The current schema version is 0.1. The pack is the
unit a collaborator should copy, validate, open in Studio, and instantiate in a
runtime. It is not merely a pointer to a machine-local URDF tree.

\begin{lstlisting}[language={},caption={Representative Robot Pack layout.}]
smores_ep/
  robot_pack.yaml
  specs/
    module_types.yaml
    connector_types.yaml
    capabilities.yaml
  mappings/
    urdf_mapping.yaml
  assets/
    urdf/smores_ep.urdf
    mujoco/smores_ep.xml
    meshes/*.stl
    collision_meshes/*.obj
\end{lstlisting}

The root manifest defines identity, semantic version, description, bounded JSON-safe
metadata, model-view recipes, asset catalogs, split-document paths, and mappings. Split
documents keep frequently edited catalogs readable and give validation issues a
specific document and JSON-pointer-like path.

\subsection{Strict schema and stable identifiers}

Pydantic models are strict, frozen, and reject unknown fields. Coercion, non-finite
numbers, and invalid paths fail rather than being normalized silently. Semantic IDs use
lowercase snake case, begin with a letter, and are at most 64 characters. This is why a
connector type ID such as \code{EP} is invalid while \code{ep\_face} is valid. Display
names remain free-form and human-readable.

The distinction matters operationally. IDs appear in cross-references, runtime
instance IDs, mappings, graph keys, and events, so casual renaming is disruptive.
Names are presentation. Studio turns raw schema messages into this distinction and
suggests a normalized ID where possible.

\subsection{Module, joint, and connector semantics}

A module type selects a mechanical asset, identifies its root link, optionally records
mass, declares semantic joints and connector instances, and advertises capabilities.
It does not replicate all URDF links. A joint records source name, joint type,
parent/child link, normalized axis, exposed control modes, and unit-bearing limits.
The schema can express position, velocity, and effort; the current \mujoco{} adapter
implements effort only.

A connector instance includes:

\begin{itemize}
  \item stable ID and connector type reference;
  \item parent link;
  \item a named frame and/or numeric local pose;
  \item docking and approach axes expressed in the parent-link frame; and
  \item bounded custom metadata.
\end{itemize}

A connector type contains the shared mating contract: active flag, gender, mutual
compatibility catalog, discrete or continuous roll orientations, capture shape and
tolerances, physical connection intent, compliance/load data, undocking support,
docking policy, and metadata. Current connection intent vocabulary includes fixed,
compliant, hinge, ball, and custom; runtime support is narrower. This is deliberate:
the schema can preserve future intent while validation and capability checks refuse
unsupported execution.

\subsection{Custom metadata}

Metadata is intentionally JSON-compatible rather than an unbounded YAML escape hatch.
It is supported at the manifest, connector type, connector instance, and model-view
configuration levels. This allows platform provenance, calibration notes, UI hints,
and external identifiers without weakening standard fields. Metadata does not alter
docking semantics unless a registered consumer interprets it. The schema presently has
no arbitrary module-type metadata field, which prevents the report from overstating
the scope of the feature.

\subsection{Backend mappings and capabilities}

Mappings associate a semantic module with a selected asset and provide link, joint,
connector-frame, and constraint name translations. The design anticipates different
naming schemes in URDF, MJCF, and a future USD backend. Current \mujoco{} support does
not yet consume the full mapping vocabulary and generally expects semantic names to
match the selected asset. Named-frame-only connectors are refused; numeric local poses
are required for the current path.

Capabilities are either primitive actions or higher-level behaviors. They can require
connector types and can be advertised by modules. In version 0.1 they are validated
descriptions, not executable command signatures. This distinction prevents a pack from
claiming that a runtime behavior engine already exists.

\subsection{Safety and loading}

The loader uses safe YAML, rejects duplicate keys, checks schema version, and injects
catalog-key IDs before model validation. Pack and split-document symlinks are rejected.
Paths must be relative, normalized, contained under the pack root, and free of home,
absolute, Windows-drive, backslash, control-character, and parent-traversal forms.
Network asset retrieval is not permitted during load.

These restrictions serve two purposes. They make a pack genuinely portable, and they
prevent a downloaded data package from reading arbitrary files. Robot Packs also name
builders and backends by stable IDs; they do not contain executable Python import
paths.

\subsection{Validation profiles}

Validation produces deterministic, structured issues with severity, stable code,
document, path, entity, explanation, and suggested fix. The \emph{authoring} profile
permits incomplete data with warnings so a draft can be opened and refined. The
\emph{simulation} profile promotes missing connector limits, joint limits, or mappings
to errors where structural readiness requires them.

Simulation validation must not be confused with a solver run. It does not parse every
referenced URDF name, load a backend, prove collision quality, or guarantee controller
stability. It states that the semantic package is complete enough to attempt a
simulation.

\subsection{Canonical writing and atomic update}

The writer revalidates the in-memory model, renders deterministic split YAML, copies
declared assets without following symlinks or special entries, stages the result,
reloads it, and compares semantics before publication. A new \code{write} refuses an
existing destination. An \code{update} builds a sibling staging directory, renames the
original to a backup, publishes the new tree, and attempts rollback if publication
fails.

Canonical output intentionally discards comments, anchors, hand formatting, and unknown
files. This is the cost of treating the schema model, rather than arbitrary YAML text,
as authoring state. Studio's Save and Export commands use the same path, so CLI and GUI
do not generate subtly different packs.

\subsection{URDF import}

The dependency-light URDF importer caps input size, forbids symlinks, DTDs, and entity
declarations, and parses robot/link/joint trees, origins, axes, limits, masses, basic
geometry, mesh references, and material colors. It resolves relative, local
\code{file://}, and \code{package://} meshes against explicit roots; it never downloads
HTTP assets. The draft builder copies resolved meshes, rewrites URDF references,
creates an initial mapping, and derives module mass.

Connectors cannot be inferred reliably from ordinary URDF. Their semantics and precise
mating frames are authored afterward. Xacro expansion, transmissions, simulator
extensions, comprehensive OBJ/DAE sidecars, and texture copying remain future work.

\subsection{Robot Pack Studio}

Studio is a native PySide6 application with a PyVista/PyVistaQt viewport. Its Qt-free
\code{StudioProject} treats each edit functionally: dump the existing pack, apply a
bounded change, revalidate, and return a dirty replacement. This centralizes
referential integrity outside individual widgets.

The current interface supports:

\begin{itemize}
  \item opening a pack or creating one from URDF;
  \item project and URDF tree navigation;
  \item pack, connector-type, connector-instance, joint, and model-view property
        editing;
  \item connector-type add/edit/reference-safe removal;
  \item connector-instance add/edit/remove and parent-link association;
  \item row-oriented JSON metadata editing;
  \item authoring and simulation-profile validation;
  \item canonical YAML preview;
  \item atomic Save and non-overwriting Export As; and
  \item a per-launch local session log truncated at startup.
\end{itemize}

The viewport renders the URDF at zero joint configuration with detailed visual meshes,
material or fallback colors, smooth/PBR shading, multiple lights, shadows, an expanded
ground grid, link frames, joint axes, connector glyphs, and optional red wireframe
collision proxies. Selecting a link can initiate connector creation, and link selection
is highlighted.

\subsection{Studio limitations}

Studio currently visualizes one module at zero configuration. It does not animate
joints, render multi-module assemblies, pick connector glyphs in 3-D, provide a
connector transform gizmo, draw capture volumes, render named-frame-only connectors,
or show live contacts and forces. Texture images and normal maps are not bundled or
rendered. Rebuilding after an edit may reset the camera, and UI-level regression
coverage is lighter than core model coverage. These limitations motivate future
integration between the authoring and runtime shells but do not change the current
authority boundary.


\section{Canonical Runtime, Docking, and Assembly Semantics}
\label{sec:runtime}

\subsection{Scene and canonical world}

A \code{SceneSpec} is an ordered collection of module placements. It validates unique
instance IDs and known module types; its grid helper creates deterministic names. At
session creation, the backend loads the scene, returns stable handles, and supplies an
initial snapshot. The runtime constructs mutable module, link, joint, connector, and
connection records under \code{WorldState}.

It is useful to qualify the term \emph{event-sourced}. Logical topology and docking
semantics are mutated through events. Backend samples directly update time, poses,
twists, joint state, measured connector frames, and transient lifecycle state. The
event log is therefore a complete record of semantic connection history, not a
lossless physical trajectory recorder.

The world exposes four monotonic revisions: backend sample sequence, topology,
docking-state, and event append count. Derived consumers can invalidate exactly what
they depend upon.

\subsection{Per-step ordering}

\begin{figure}[htbp]
\centering
\begin{tikzpicture}[
  node distance=5mm,
  step/.style={draw,rounded corners,fill=modsimblue!10,draw=modsimblue,
               minimum width=0.78\textwidth,minimum height=8mm,align=center},
  arrow/.style={-{Latex[length=2mm]},thick}
]
  \node[step] (s1) {1. Backend advances physics by $\Delta t$};
  \node[step,below=of s1] (s2) {2. Backend snapshot reports link, joint, connector, and force observations};
  \node[step,below=of s2] (s3) {3. World ingests the fresh sample and resolves connector frames/velocities};
  \node[step,below=of s3] (s4) {4. Overload evaluation requests releases};
  \node[step,below=of s4] (s5) {5. Docking processes releases, candidates, guards, and commits};
  \node[step,below=of s5] (s6) {6. Revisions, events, metrics, and derived views become observable};
  \foreach \a/\b in {s1/s2,s2/s3,s3/s4,s4/s5,s5/s6} \draw[arrow] (\a) -- (\b);
\end{tikzpicture}
\caption{One canonical runtime step. Docking decisions use state reported after the
backend step, and releases are processed before new commits.}
\label{fig:step-pipeline}
\end{figure}

This order prevents decisions against stale geometry. Processing releases before
detection also prevents a connection from breaking and reforming before its release is
visible. A redock cooldown should be authored when two just-released surfaces remain in
capture range.

\subsection{Connector frame and point velocity}

For a connector $C_i$ attached to link $L_i$, an authored numeric pose gives

\begin{equation}
{}^{W}\mathbf T_{C_i} =
{}^{W}\mathbf T_{L_i}\,{}^{L_i}\mathbf T_{C_i}.
\label{eq:connector-transform}
\end{equation}

If the backend materializes and reports the connector frame, that measurement takes
precedence over recomposition. This keeps acceptance consistent with engine geometry,
especially for connectors on articulated child links.

The point velocity includes angular motion:

\begin{equation}
\mathbf v_{C_i} = \mathbf v_{L_i} +
\boldsymbol\omega_{L_i}\times(\mathbf p_{C_i}-\mathbf p_{L_i}).
\label{eq:connector-velocity}
\end{equation}

Without the lever-arm term, a rotating connector could be classified as stationary and
dock above its declared speed limit.

\subsection{Broad phase and deterministic candidates}

Free resolved connectors are indexed in a uniform spatial hash whose cell size is the
largest position tolerance in the pack, with a small positive lower bound. Only nearby
cells produce opportunistic pairs. Explicitly commanded pairs bypass broad-phase
distance rejection so an invalid command produces an informative failure rather than
silence. Pairs are sorted before evaluation.

\subsection{Compatibility}

Type compatibility must be mutual. If type $A$ names $B$ but $B$ does not name $A$,
the pair is rejected as an authoring error rather than interpreted asymmetrically.
Genders pair male--female, hermaphroditic--hermaphroditic, and
genderless--genderless. A connector cannot dock to itself or another connector on the
same module.

\subsection{Geometric and dynamic acceptance}

Acceptance produces a structured result instead of a boolean. For each tolerance $x$,
the effective pair limit is the stricter endpoint declaration,

\begin{equation}
\epsilon_x^{ij}=\min(\epsilon_x^i,\epsilon_x^j).
\label{eq:tolerance}
\end{equation}

Both connector-local capture shapes must contain the relative offset. A sphere uses
Euclidean distance; a box bounds each local coordinate; a cylinder separately bounds
axial and radial displacement.

Two docking axes mate antiparallel, so the axis error is

\begin{equation}
e_a=\cos^{-1}\!\left(\hat{\mathbf a}_i^\top(-\hat{\mathbf a}_j)\right).
\label{eq:axis-error}
\end{equation}

Roll is measured about the shared docking axis using projected connector-frame
reference directions. For a discrete allowed set $\Phi$, the governing error is

\begin{equation}
e_\phi=\min_{\phi_k\in\Phi}
\left|\operatorname{wrap}(\phi-\phi_k)\right|.
\label{eq:roll-error}
\end{equation}

Both endpoints are checked with opposite roll signs. The committed connection stores
the matched orientation index because configuration recognition cannot recover it
reliably from a later noisy pose. Relative connector-point speed is the fourth
criterion. The result reports each measurement, tolerance, margin, and failure reason.

\subsection{Guards and docking policy}

Guards operate after geometry. They reject engaged connectors, cooldown violations,
undefined or conflicting physical intent, unsupported connection types, and missing
commands when the pair is not jointly auto-latching. Closing a loop within an existing
assembly is allowed but warned because rings and lattices are legitimate while
redundant rigid constraints can be accidental.

Alignment policy is either \emph{measured} or \emph{nominal}. Measured alignment freezes
the observed relative pose. Nominal alignment asks the backend to snap connector frames
to their ideal antiparallel/roll pose. If either endpoint requests nominal alignment,
the pair snaps. Physical controller experiments generally prefer measured alignment;
lattice experiments may prefer nominal alignment to avoid accumulated drift.

\subsection{Two-phase physical/logical commit}

\begin{figure}[htbp]
\centering
\begin{tikzpicture}[
  node distance=7mm and 12mm,
  box/.style={draw,rounded corners,align=center,minimum width=35mm,minimum height=9mm},
  decision/.style={diamond,draw,aspect=2.2,align=center,inner sep=1.5pt},
  arrow/.style={-{Latex[length=2mm]},thick}
]
  \node[box,fill=modsimblue!10] (eval) {Recheck acceptance\\and guards};
  \node[box,fill=modsimgold!20,right=of eval] (request) {Create physical\\connection request};
  \node[decision,right=of request] (ok) {Backend\\accepted?};
  \node[box,fill=modsimcyan!15,above right=6mm and 12mm of ok] (commit)
    {Append DockCommitted\\then AssemblyMerged};
  \node[box,fill=modsimorange!12,below right=6mm and 12mm of ok] (fail)
    {Append DockFailed\\with reason/detail};
  \draw[arrow] (eval) -- (request);
  \draw[arrow] (request) -- (ok);
  \draw[arrow] (ok) -- node[above,sloped,font=\scriptsize]{yes} (commit);
  \draw[arrow] (ok) -- node[below,sloped,font=\scriptsize]{no} (fail);
\end{tikzpicture}
\caption{Two-phase connection commit. Canonical topology changes only after the
backend confirms the physical constraint.}
\label{fig:two-phase}
\end{figure}

The backend is first asked to create a constraint. Only a successful returned handle
permits \code{DockCommitted}. If the endpoints previously belonged to separate
components, \code{AssemblyMerged} follows. Backend refusal, unsupported capabilities,
or resource exhaustion records \code{DockFailed}; core never fabricates a connection.

Undocking is symmetric: remove the backend constraint first, then append
\code{UndockCommitted}, then recompute the affected component and append
\code{AssemblySplit} if needed. A permanent connector refuses undocking.

\subsection{Assemblies and topology}

At time $t$, let

\begin{equation}
G_t=(V,E_t), \qquad \mathcal A_t=\operatorname{CC}(G_t),
\label{eq:assemblies}
\end{equation}

where $V$ is the fixed set of module instances and $E_t$ is the set of committed
connections. Parallel edges are preserved even though connected-component membership
does not require multiplicity. Merge joins indexed component sets. Release recomputes
only the affected component rather than the entire world. Assembly identity derives
from membership so it is stable across processes and replay.

\subsection{Events and metrics}

The append-only event vocabulary includes candidate detection, dock success/failure,
undock success/failure, assembly merge/split, and connector overload. Appends receive a
monotonic zero-based sequence. Opportunistic near-miss failures are not logged every
step because they would overwhelm history; explicit failed commands are logged.

Metrics are computed from canonical state and history: module/free/connected counts,
assembly count and maximum size, connector lifecycle counts, active connections,
docking and undocking successes/failures/attempts, merges, splits, overloads, and event
count.

\subsection{Joint commands}

\code{JointCommand} carries a stable module/joint ID, semantic control mode, and scalar
SI target. A batch is validated fully before any backend mutation: IDs must be unique
and known, the pack must declare the mode and applicable limits, the backend must
support the mode, and values must be finite non-booleans. A failure rejects the entire
batch. Commands persist in a backend until replaced or cleared, which makes a complete
batch important for preventing stale effort after a topology change.

\subsection{Overload boundary}

If a backend reports constraint-force magnitude and the two connector policies declare
break-force thresholds, the weaker threshold governs. Exceeding it emits
\code{ConnectorOverloaded} and initiates release. The mock can inject forces for tests;
the current \mujoco{} adapter does not yet report equality force, so physical breakage
is not active there. Normal/shear/bending values stored under connector limits remain
descriptive until a richer load model is implemented.


\section{Backend Contract and MuJoCo Physics Realization}
\label{sec:backends}

\subsection{Adapter contract}

The base adapter contract contains seven responsibilities: advertise capabilities,
load a pack and scene, step, snapshot, create a physical connection, remove it, and
shutdown. The snapshot is intentionally small: time, link poses and twists, scalar
joint state, optional measured connector frames, and optional constraint-force
magnitudes. Contacts, solver internals, renderer objects, and engine arrays do not leak
into canonical core state.

Optional structural protocols add joint commands and direct module kinematics. Direct
pose/twist/wrench methods are useful for deterministic scene staging and kinematic
fixtures, but physical demonstrations enforce that these are not called after initial
placement.

\subsection{Lazy registration and selection}

Backends register a name, summary, loader, and installation hint. Resolution is an
explicit argument, then \code{MODSIM\_BACKEND}, then \code{mock}. Loading is lazy, so a
base installation can import the entire core without importing \mujoco{}. A backend
implementation can live in a separate distribution and register through the same
mechanism.

\subsection{Mock backend}

The mock backend is kinematic and dependency-free. It has no gravity, contact, or
articulated link motion. It nonetheless implements the semantic behaviors core needs:
module placement, world-frame twist integration, physical-connection success/failure,
rigid propagation of connected groups, release, injected loads, and snapshot
generation. It is therefore a reference semantics backend, not a low-fidelity physics
claim.

\subsection{MuJoCo scene composition}

The \mujoco{} adapter composes all module placements into one model. For each instance
it:

\begin{enumerate}
  \item selects a same-ID MJCF asset when available, otherwise imports URDF;
  \item attaches the module under a namespace prefix \code{<module>/};
  \item adds a free joint to the module root;
  \item creates connector sites from authored local poses;
  \item creates bounded unit-gear motors for declared effort-controlled scalar joints;
  \item reserves inactive weld and contact-exclusion slots; and
  \item compiles and indexes bodies, joints, sites, actuators, and assets.
\end{enumerate}

The default gravity vector is $(0,0,-9.81)\,\mathrm{m/s^2}$. An optional engine plane is
physically infinite; its visual half-extent is ten meters. Visual, environment, and
debug collision geoms use distinct groups. When importing URDF, the adapter retains
visual geometry unless the asset explicitly requests otherwise. Collision geoms are
placed in hidden group 3 when separate visuals exist, so physics remains active without
opaque proxies obscuring the robot.

\subsection{Time stepping and snapshots}

The backend covers a requested duration with rounded whole solver steps and reports the
actual engine time. Snapshot extraction resolves world body motion, connector site
frames, joint position/velocity, and actuator effort. For anisotropic \smoresep{} tire
contacts, stepping is split so the contact tangent basis can be oriented with the
current wheel axle before constraint solution.

\subsection{Runtime weld pool}

\mujoco{} model topology is compiled, so a runtime dock claims a preallocated inactive
weld rather than adding a new equality object. Each slot has a paired precompiled body
contact exclusion. Claiming a connection sets the equality body operands, relative
transform data, and active flag; it publishes a sorted exclusion signature so the
welded body pair does not simultaneously fight a contact constraint. Release
deactivates and returns both resources.

The connector-relative request must be converted to a body-relative weld:

\begin{equation}
\mathbf T_{B_aB_b}=
\mathbf T_{B_aC_a}
\mathbf T_{C_aC_b}
\mathbf T_{C_bB_b}.
\label{eq:weld-transform}
\end{equation}

The implementation isolates this transformation in a testable module. Nominal snapping
also accounts for connectors on articulated child bodies: it measures root-to-child
transform, solves the desired root pose, clears root and articulated velocity, and only
then activates the weld.

Pool exhaustion and unsupported connection types are ordinary backend refusals. They
propagate to \code{DockFailed} rather than a logical-only success. The current adapter
supports fixed connections only. Compliance, hinge, ball, and custom constraints remain
future work.

\subsection{Effort actuation and state}

The adapter advertises the effort control mode. Scene composition creates one actuator
for each semantically declared effort-controlled joint with a finite effort limit.
Atomic commands resolve all IDs and bounds before writing \code{data.ctrl}. Values
persist until replaced or cleared. Snapshots report position, velocity, and actuator
effort using stable semantic joint IDs.

The architecture can represent position and velocity modes, but the adapter does not
yet implement them. Likewise, the current system does not infer arbitrary transmissions
from URDF. A future actuator catalog should make motor, gear, and controller dynamics
explicit.

\subsection{Native viewer and pacing}

The native passive viewer starts with collision debug group 3 hidden and preserves the
final state for inspection. On macOS, \mujoco{} requires \code{mjpython} for passive
viewer main-thread ownership \citep{mujocoPythonDocs}; the Runtime Inspector resolves a
sibling environment executable automatically.

The launch-time real-time factor $\rho>0$ changes only a wall-clock deadline:

\begin{equation}
t_w^{\mathrm{deadline}} = t_{w,0}+\frac{\Delta t_s}{\rho}.
\label{eq:pacing}
\end{equation}

Simulation timestep, number of steps, motor targets, event timestamps, and simulated
duration remain unchanged. Values below one create slow motion; values above one request
accelerated playback. If physics and rendering cannot keep up, execution proceeds as
fast as possible without skipping physics steps.

At the current head commit, passive viewer synchronization is capped at approximately
30 Hz rather than occurring at every solver step. Synchronization is scheduled from
completion so an expensive render does not trigger another render on the immediately
following physics step, and a final synchronization is guaranteed when the run ends
between refreshes. This decoupling is important at a \SI{2}{ms}-style solver step,
where sync-per-step would attempt 500 visual refreshes per simulated second.

\subsection{Backend comparison}

\begin{table}[htbp]
\centering
\caption{Implemented backend capabilities.}
\label{tab:backend-comparison}
\begin{tabularx}{\textwidth}{L{0.30\textwidth}YY}
\toprule
Capability & Mock & MuJoCo \\
\midrule
Dependencies & Base only & Optional \code{mujoco} and NumPy \\
Dynamics/contact/gravity & No & Yes \\
Articulated kinematics & No & Yes \\
Joint feedback & No & Position, velocity, effort \\
Joint command & No & Effort \\
Measured connector frames & Composed & Connector sites \\
Runtime fixed connection & Yes & Preallocated weld \\
Runtime release & Yes & Weld deactivation \\
Paired contact exclusion & Not applicable & Yes \\
Constraint force reporting & Injectable test value & Not yet \\
Native 3-D viewer & No & Yes \\
\bottomrule
\end{tabularx}
\end{table}

\subsection{Conformance}

Cross-backend tests run a common pack and scenario with gravity disabled and compare
semantic modules, assemblies, connector frames, connection identities, mating
orientation, and committed event sequence. This isolates adapter correctness from
expected physical trajectory differences. The mock is not evidence of a second
physical solver, but it is a valuable executable specification of the semantic
contract.


\section{Model Views, Runtime Inspector, and Human Observability}
\label{sec:views}

\subsection{Why topology is derived}

A module graph is an excellent abstraction for reconfiguration algorithms, but it
cannot be canonical physical state. It omits links, joint values, connector frames,
contact, pose, and transient docking lifecycle. Making it independently mutable would
create the duplicated-truth problem \modsim{} is designed to avoid.

Model views are therefore strict, immutable, JSON-safe Pydantic data-transfer objects.
They contain a schema version, source stamp, simulation time, pack identity, and typed
content. They never contain Qt objects, \mujoco{} handles, NumPy arrays, NetworkX
graphs, or references to mutable \code{WorldState}. A client may convert the DTO into
its chosen rendering or analysis library.

\subsection{Factory and builder abstraction}

\code{ModelViewBuilder} is the extension point. A builder declares a stable ID, name,
result type, supported modes, whether it requires a live world, a dependency token,
and a build method. \code{ModelViewFactory} is per-client rather than global. It
registers built-ins, rejects duplicate IDs, lists descriptors, resolves pack recipes,
builds one or all views, and maintains a bounded thread-safe cache.

Robot Pack recipes contain only builder ID, modes, default flag, display name, and
bounded JSON configuration. They contain no executable import path. An unknown builder
remains portable valid data but cannot be constructed until an application registers an
implementation.

\subsection{Module topology graph}

The first built-in view is a runtime multigraph:

\begin{itemize}
  \item every module instance is a node, including isolated modules;
  \item every committed connection is one undirected edge;
  \item parallel connections remain separate edges;
  \item nodes copy module pose and assembly identity;
  \item edges copy endpoints, connection ID, orientation index, and commit time; and
  \item deterministic sorting makes serialization stable.
\end{itemize}

Only \code{DockCommitted} and \code{UndockCommitted} change graph topology. Candidates,
failed attempts, and planned edges are not drawn as committed connections.

\subsection{Revision-aware caching}

The topology builder observes the backend sample sequence because node poses can move,
and the topology revision because edges can change. A cache entry is keyed by source
identity, recipe fingerprint, and builder token. Unrelated event or docking revisions
can update a source stamp without requiring a full graph reconstruction. The cache is
least-recently-used and bounded, retaining the newest views rather than growing with
simulation length.

\subsection{Runtime Inspector frame}

After stepping, the runtime owner builds a \code{RuntimeInspectorFrame} containing:

\begin{itemize}
  \item the immutable graph view;
  \item scenario phase, plan/action index, and human-readable detail;
  \item current runtime metrics;
  \item a contiguous event delta since the last successfully published frame; and
  \item backend, session, time, and source revision information.
\end{itemize}

Graph frames may be refresh-limited, but event deltas are lossless and contiguous. The
event cursor advances only after a frame is successfully constructed for publication.

\subsection{Two execution lanes}

The Runtime Inspector supports two ownership patterns that share the same frame
contract.

\paragraph{Semantic-only lane.}
With \code{--no-viewer}, a dedicated QThread worker owns the runtime, scenario,
backend, and model-view factory. It paces against simulated time and emits immutable
frames to the Qt thread at a configured wall-clock publication rate.

\paragraph{Native-viewer lane.}
With \mujoco{} viewer enabled, a companion process owns the one authoritative runtime
and native viewer. This is required on macOS because Qt and the passive \mujoco{} viewer
cannot both own the Cocoa main GUI thread. The Qt parent displays topology/events; the
child advances physics and displays geometry. They are not two synchronized
simulations.

\begin{figure}[htbp]
\centering
\begin{tikzpicture}[
  node distance=7mm and 18mm,
  proc/.style={draw,rounded corners,minimum width=0.39\textwidth,minimum height=18mm,
               align=center,fill=modsimgray},
  item/.style={draw,rounded corners,align=center,font=\small,minimum width=0.31\textwidth,
               minimum height=8mm},
  arrow/.style={-{Latex[length=2mm]},thick}
]
  \node[proc,fill=modsimblue!10,draw=modsimblue] (qt) {Qt parent process\\Runtime Inspector};
  \node[proc,fill=modsimcyan!10,draw=modsimcyan!80!black,right=of qt] (host)
    {MuJoCo child process\\authoritative runtime + native viewer};
  \node[item,below=of qt] (present) {Presenter, stable graph layout, event table};
  \node[item,below=of host] (owner) {RuntimeSession, scenario, ModelViewFactory};
  \node[item,below=of owner] (engine) {MjModel / MjData / passive viewer};
  \draw[arrow] (host) -- node[above,font=\scriptsize]{versioned immutable frames} (qt);
  \draw[arrow] (qt) -- node[below,font=\scriptsize]{initialize / stop} (host);
  \draw[arrow] (qt) -- (present);
  \draw[arrow] (host) -- (owner);
  \draw[arrow] (owner) -- (engine);
\end{tikzpicture}
\caption{Dual-window Runtime Inspector architecture. The protocol carries values only;
the Qt process never traverses mutable world or engine objects.}
\label{fig:runtime-processes}
\end{figure}

\subsection{Versioned local protocol}

The child emits newline-delimited strict JSON records prefixed by
\code{MODSIM\_RUNTIME/1}. Message kinds are hello, initialize, stop, status, frame,
error, and finished. The framer handles fragmented reads and rejects bad prefixes,
UTF-8, multiline documents, oversized records, and unterminated input. Standard output
is reserved for protocol data; standard error carries diagnostics that the parent
mirrors into its session log.

Shutdown is cooperative: the parent sends Stop, waits, then escalates to terminate and
kill only after deadlines. End-of-file also stops the child. The child exits the viewer,
shuts down the backend, sends a terminal record, and joins its input watcher. These
details prevent orphan viewers and interpreter aborts from blocked standard-input
threads.

\subsection{Presentation behavior}

The Qt-free presenter owns layout and selection, not runtime state. It assigns stable
deterministic node positions, keeps nodes in place across pose samples, renders parallel
edges as distinct curves, removes an edge selection when that edge disappears, and
retains node selection. It rejects regressing source stamps and gaps or conflicts in
event deltas.

The window shows graph, event table, scenario phase/detail, backend, simulation time,
module/assembly/connection counts, dock/undock outcomes, event count, and source
revisions. A target-speed label explicitly states the requested pacing factor. The
native viewer renders physical geometry and contact; the semantic window does not.

\subsection{Current visualization limits}

Only the module-topology graph is rendered. There are no frame, port, matrix, lattice,
or hybrid-docking panels; no time-series plots; no pause, restart, or scrub; no manual
joint controls; and no general scene editor. The Runtime Inspector remains a separate
window from Robot Pack Studio. These are presentation limitations rather than reasons
to weaken the immutable view boundary.


\section{SMORES-EP Case Study: Model, Contact, and Control}
\label{sec:smores}

\subsection{Platform basis}

\smoresep{} was selected because it exercises nearly every architectural concern in a
compact platform: reusable articulated modules, four connector faces, differential
drive, active reconfiguration, discrete mating orientations, and published connector
and topology research. The official platform description identifies four active DOFs
and rubber-tired left/right wheels \citep{modlabSmoresEp}. The EP-Face is genderless and
switchable, with published holding-force and capture characterization
\citep{tosun2016epface}. Self-assembly work divides docking into navigation, pose
adjustment, and final approach \citep{liu2020parallel,liu2023smoresep}.

The committed pack is a research artifact, not a complete hardware specification. This
section therefore labels parameter provenance. Superscripts mean published
(\sourcefact{}), geometry-derived (\cadfact{}), provisional simulation choice
(\simchoice{}), or measured/tested simulation outcome (\measuredfact{}).

\subsection{Pack contents}

The \pathcode{examples/robot_packs/smores_ep} pack contains:

\begin{itemize}
  \item one module type and one genderless self-compatible connector type;
  \item four connector instances named \code{bottom}, \code{pan}, \code{left}, and
        \code{right};
  \item dock and undock capability descriptions;
  \item a default runtime topology view recipe;
  \item a Fusion-derived URDF and a same-ID physics-specific MJCF asset;
  \item five detailed STL visual meshes and simplified collision/contact assets; and
  \item a URDF backend mapping.
\end{itemize}

The module mass is \SI{0.4386368}{kg}\cadfact{}. LEFT and RIGHT wheel joints and PAN are
continuous; TILT is bounded to $\pm\pi/2$. The wheel speed cap is
$\pi/2\,\mathrm{rad/s}$ (\SI{90}{\degree\per\second})\sourcefact{}. The wheel effort limit
\SI{0.04}{N.m}\simchoice{} and PAN/TILT limits \SI{0.10}{N.m}\simchoice{} are
simulation bootstrap values, not identified motor constants.

\subsection{Connector frames}

\begin{table}[htbp]
\centering
\caption{Committed \smoresep{} connector instances. Positions are in parent-link
coordinates.}
\label{tab:smores-connectors}
\begin{tabularx}{\textwidth}{L{0.12\textwidth}L{0.19\textwidth}L{0.28\textwidth}Y}
\toprule
ID & Parent link & Position (m) and RPY (rad) & Docking axis \\
\midrule
\code{bottom} & \code{base\_link} & $(-0.0107416,0,0)$; $(0,0,\pi)$ & $(-1,0,0)$ \\
\code{pan} & \code{front\_wheel\_1} & $(0.005762,0,0)$; $(0,0,0)$ & $(1,0,0)$ \\
\code{left} & \code{left\_wheel\_1} & $(0,0.04405,0)$; $(0,0,\pi/2)$ & $(0,1,0)$ \\
\code{right} & \code{right\_wheel\_1} & $(0,-0.04405,0)$; $(0,0,-\pi/2)$ & $(0,-1,0)$ \\
\bottomrule
\end{tabularx}
\end{table}

The corrected BOTTOM connector is on the rigid rear $-X$ face of the base, not the
underside. Its plane is derived from the dominant rear surface in the Fusion mesh. This
correction is essential for TOP-to-BOTTOM snake connections: an underside frame turns a
nose-to-tail chain into an unintended vertical or kinked structure.

\subsection{EP-face semantic model}

The connector type is genderless and self-compatible. Allowed roll states are
$\{0,\pi/2,\pi,3\pi/2\}$\sourcefact{}. The current capture proxy is cylindrical with
\SI{6}{mm} position tolerance\simchoice{}, \SI{25}{\degree} orientation tolerance, and
\SI{0.05}{m/s} maximum relative speed. The 25-degree value is motivated by published
residual attraction under angular misalignment \citep{tosun2016epface}, but it should
not be read as a calibrated complete capture model.

Connection intent is fixed, undocking is supported, auto-latch is disabled, measured
alignment is selected, and redock cooldown is \SI{0.1}{s}. Listed normal, shear, and
bending loads are explicitly unverified provisional metadata. The current generic
acceptance model uses one radial position bound; it does not reproduce the literature's
separate normal and lateral capture envelopes.

\subsection{Visual and contact models}

Detailed STLs carry visible shape but no textures or material graph. The URDF assigns a
uniform silver material. In MJCF, visual meshes have zero density and do not carry
contact. Contact uses simpler primitives:

\begin{itemize}
  \item two tire cylinders with radius \SI{0.040}{m}\cadfact{} and half-width
        \SI{0.01045}{m}\cadfact{};
  \item axle track \SI{0.0672}{m}\cadfact{};
  \item a small rear support skid\simchoice{} to make the two-wheel chassis statically
        supportable; and
  \item connector-face collision proxies in a separate category, so faces can meet
        without dragging against the ground.
\end{itemize}

Tire/ground contact uses \code{condim=4} and friction tuple
$(0.05,1.0,0.01,10^{-4},10^{-4})$\simchoice{}. The low/high tangential values establish
anisotropy; the backend rotates the tangent frame with the wheel axle. Skid/ground
friction is $(0.08,0.08,0.005,10^{-4},10^{-4})$\simchoice{}. Wheel damping and armature
are \num{0.001} and \num{0.0005}; PAN/TILT damping and armature are \num{0.01} and
\num{0.0002}\simchoice{}. These quantities were tuned for stable first-pass dynamics and
must eventually be identified from hardware.

\subsection{Differential-drive kinematics}

For wheel radius $r$, track $b$, and left/right angular speeds
$\omega_L,\omega_R$, the ideal planar body twist is

\begin{align}
v &= \frac{r}{2}(\omega_R+\omega_L), \\
\dot\psi &= \frac{r}{b}(\omega_R-\omega_L).
\label{eq:diff-drive}
\end{align}

The inverse targets are

\begin{align}
\omega_R^\star &= \frac{v+\frac{b}{2}\dot\psi}{r}, &
\omega_L^\star &= \frac{v-\frac{b}{2}\dot\psi}{r}.
\label{eq:wheel-targets}
\end{align}

Each wheel uses bounded proportional velocity-to-effort feedback,

\begin{equation}
\tau_i = \operatorname{sat}_{\tau_{\max}}
\left[k_v(\omega_i^\star-\omega_i)\right],
\label{eq:wheel-effort}
\end{equation}

with $k_v=0.1\,\mathrm{N,m/(rad/s)}$\simchoice{} and
$\tau_{\max}=0.04\,\mathrm{N,m}$\simchoice{}. PAN and TILT use bounded PD hold,

\begin{equation}
\tau_h = \operatorname{sat}_{\tau_{h,\max}}
\left[k_p(q^\star-q)-k_d\dot q\right],
\label{eq:hold-effort}
\end{equation}

with $k_p=1.0$, $k_d=0.02$, and \SI{0.1}{N.m} cap\simchoice{}. Current feedback runs at
the solver period for stability. A calibrated actuator model should permit a more
hardware-faithful outer-loop frequency.

\subsection{Rigid-assembly wheel allocation}

The final two Driver-to-Snake actions move a connected three-module component. A
desired planar twist is evaluated at every module center. For module $i$ at position
$\mathbf p_i$, reference point $\mathbf p_0$, reference linear velocity $\mathbf v_0$,
and desired yaw rate $\omega$, the rigid velocity target is

\begin{equation}
\mathbf v_i^\star = \mathbf v_0 +
\omega\,\hat{\mathbf z}\times(\mathbf p_i-\mathbf p_0).
\label{eq:rigid-velocity}
\end{equation}

The allocator projects $\mathbf v_i^\star$ onto each module's planar forward axis and
uses \cref{eq:wheel-targets}. The orthogonal projection is a nonholonomic residual. If
the largest requested residual exceeds \SI{0.08}{m/s}\simchoice{}, the whole twist is
scaled. Wheel targets are then scaled again if any exceeds $\pi/2\,\mathrm{rad/s}$.
Both requested and applied residual maxima are retained for diagnosis.

This is a pragmatic controller for a rigid assembly whose module wheel axes are not
all located at the reference center. It is not an optimal traction allocator, and the
reported residual is a command-space diagnostic, not measured tire slip.

\subsection{Two-module physical scenario}

The first physical scenario uses a stationary module's rear BOTTOM connector and a
moving module's PAN/TOP connector. Creation performs all capability, joint-feedback,
and effort-command preflight before staging the pair once. It then:

\begin{enumerate}
  \item settles both free modules for \SI{0.5}{s};
  \item brakes the target and drives the mover forward through wheel effort;
  \item steers from measured connector geometry and continues until a strict
        \SI{1}{mm} near-contact gate is met;
  \item requests the ordinary measured-frame fixed connection;
  \item holds for \SI{1}{s};
  \item waits \SI{80}{ms}\sourcefact{} as a modeled release delay;
  \item removes the weld; and
  \item reverses \SI{0.04}{m} through wheel effort.
\end{enumerate}

No root pose, root twist, or external wrench is written after initial staging. The
approach and separation therefore depend on gravity, contact, inertia, joint motion,
and effort limits. Magnetic attraction before capture is absent; the accepted latch is
an ideal fixed weld.

\subsection{Driver-to-Snake plan}

The seven-module plan follows the connector substitutions reported by
\citet{liu2019distributed}. The initial Driver tree and replacements are listed in
\cref{tab:snake-actions}. The final physical module order is
$1-3-2-4-5-6-7$.

\begin{table}[htbp]
\centering
\caption{Driver-to-Snake connector replacements. TOP maps to \code{pan}; BOTTOM maps to
\code{bottom}.}
\label{tab:snake-actions}
\begin{tabularx}{\textwidth}{L{0.08\textwidth}YY}
\toprule
Action & Release & Replacement \\
\midrule
1 & \code{module\_1/bottom}--\code{module\_2/pan} &
    \code{module\_1/pan}--\code{module\_3/bottom} \\
2 & \code{module\_7/pan}--\code{module\_5/bottom} &
    \code{module\_7/bottom}--\code{module\_6/pan} \\
3 & \code{module\_2/right}--\code{module\_4/left} &
    \code{module\_2/pan}--\code{module\_4/bottom} \\
4 & \code{module\_5/left}--\code{module\_4/right} &
    \code{module\_5/bottom}--\code{module\_4/pan} \\
\bottomrule
\end{tabularx}
\end{table}

Actions 1 and 2 move single leaves. After they dock, action 3 moves the connected
$\{1,3,2\}$ component, and action 4 moves $\{5,6,7\}$. The plan is separated from both
the kinematic and physical scenarios.

\subsection{Authored physical routes}

The physics realization stages the initial six-edge tree upright at time zero. After
that boundary, it permits only joint effort and ordinary dock/undock requests. Each
action waits through the release delay, removes the old weld, follows target-relative
clearance waypoints, performs a slow measured approach, and commits after all connector
criteria pass. Routes are recorded under \pathcode{examples/scenarios}; they are not
hidden in the backend.

The routes include empirical details such as a \SI{8}{mm} lateral prebias for the final
reverse-moving three-module train and a broad first clearance tolerance. These are
transparent tuning choices. They demonstrate the framework's physical execution path,
not robustness to arbitrary initial conditions or autonomous obstacle avoidance.

\subsection{Physics ownership invariant}

The end-to-end regression replaces backend root pose, twist, and wrench methods with
functions that fail if called after scenario creation. Completing all four replacements
under that guard is stronger evidence than merely observing wheel rotation: it proves
that the scenario cannot secretly teleport its roots after staging.

\section{Experiments, Demonstrations, and Results}
\label{sec:experiments}

\subsection{Evaluation questions}

The current evidence is organized around five evaluation questions.

\begin{description}[style=nextline]
  \item[EQ1: Package correctness]
  Do strict load, validation, mutation, persistence, and import operations reject
  malformed or unsafe data while preserving valid semantics?

  \item[EQ2: Semantic runtime correctness]
  Do docking, undocking, assemblies, events, metrics, and graph views remain mutually
  consistent under success and failure?

  \item[EQ3: Backend realization]
  Does the same semantic contract execute through the mock and \mujoco{} boundaries,
  including physical resource refusal and runtime constraint release?

  \item[EQ4: Physical actuation]
  Can articulated wheel effort, gravity, tire contact, measured connector frames, and
  runtime welds complete a dock/undock cycle without root motion injection?

  \item[EQ5: Multi-module reconfiguration]
  Can the physics path execute all four Driver-to-Snake connector replacements,
  including motion of connected three-module components, while preserving expected
  topology and bounded state?
\end{description}

The implementation does not yet answer statistical robustness, hardware fidelity, or
large-scale performance questions. Those are explicitly deferred to
\cref{sec:future}.

\subsection{Verification environment and reproducibility}

The current report snapshot is \modsim{} 0.1.0 at commit \code{4a5ca00}. A full
repository test run under Python 3.12 with the Studio and \mujoco{} dependencies
installed collected and passed 537 tests in 166.91 seconds. Branch-aware coverage of
the backend-neutral \code{modsim} package was 86\%. The wall time is reported only as
the cost of this verification run; it is not a controlled physics benchmark because
machine load, GUI libraries, coverage instrumentation, and test composition affect it.

Both committed example packs passed fresh simulation-profile structural validation
with zero errors, warnings, or informational issues:

\begin{lstlisting}[language=bash]
modsim pack validate examples/robot_packs/generic_cube --profile simulation
modsim pack validate examples/robot_packs/smores_ep --profile simulation
\end{lstlisting}

The test suite includes core unit tests, loader and writer fault cases, schema and
validator cases, mock/\mujoco{} conformance, Studio project and logging tests,
Runtime Inspector protocol/process/widget tests, and physical end-to-end tests. CI
separates core Python 3.11/3.12, headless \mujoco{}, Xvfb Studio, and packaging jobs.

\subsection{Experiment matrix}

\begin{table}[htbp]
\centering
\caption{Named demonstration matrix. Durations and steps are simulated quantities.}
\label{tab:demo-matrix}
\begin{tabularx}{\textwidth}{L{0.24\textwidth}L{0.19\textwidth}L{0.10\textwidth}L{0.18\textwidth}Y}
\toprule
Demo & Motion model & Modules & Duration / step & Expected edge history \\
\midrule
\code{dock} & kinematic approach & 2 & 4 s / 0.002 s & $0\rightarrow1$ \\
\code{dock\_undock} & kinematic approach/retract & 2 & 6 s / 0.002 s & $0\rightarrow1\rightarrow0$ \\
\code{smores\_diff\_drive\_dock\_undock} & wheel/contact physics & 2 & 12 s / 0.002 s & $0\rightarrow1\rightarrow0$ \\
\code{smores\_driver\_to\_snake} & kinematic component staging & 7 & 14 s / 0.002 s & four $6\rightarrow5\rightarrow6$ \\
\code{smores\_physical\_driver\_to\_snake} & wheel/contact physics & 7 & 210 s / 0.002 s & four $6\rightarrow5\rightarrow6$ \\
\bottomrule
\end{tabularx}
\end{table}

\subsection{Robot Pack and authoring experiments}

Schema tests cover strict scalar types, ID grammar, unique catalogs, connector
orientation sets, policies, metadata, model-view recipes, and round-trip
serialization. Loader tests cover duplicate keys, version mismatch, split-document
issues, missing files, traversal, and symlinks. Writer tests inject failures during
staging and publication to check rollback and confirm that a reloaded pack is
semantically equal to the source.

URDF tests cover geometry, origins, axes, limits, multiple visuals/collisions, material
colors, package and local mesh resolution, draft asset rebasing, unsafe XML constructs,
and unresolved assets. Studio project tests verify functional mutation, referential
connector-type removal, connector parent-link validation, metadata persistence,
mapping cleanup, model-view CRUD, Save, reopen, and Export.

The committed \smoresep{} test loads the real asset tree, verifies zero validation
issues, confirms all four connector IDs and the corrected rear BOTTOM geometry, and
optionally compiles and docks two real modules through \mujoco{}.

\subsection{Docking semantic experiments}

The docking suite exercises mutual compatibility, all gender combinations, capture
shapes, antiparallel axes, discrete roll, tighter pair tolerances, connector-point
velocity, command policy, cooldown, loop closure, unsupported physical intent,
backend refusal, deterministic arbitration, release, overload, and assembly evolution.

Event expectations are exact rather than aggregate. A successful first connection
records candidate detection, dock commit, and assembly merge. A successful release
records undock commit and assembly split. Tests assert that a backend failure cannot
create a graph edge, and conformance tests compare committed decisions across installed
backends with gravity disabled.

\subsection{Two-module physical docking experiment}

The physical pair begins with two upright modules separated by a \SI{20}{mm} connector
gap. Gravity and the ground are enabled. After settling, the moving module drives its
PAN/TOP face into the target's rear BOTTOM face. The physical regression forbids any
post-creation root pose, twist, or wrench call and asserts actual wheel-joint travel.

\begin{table}[htbp]
\centering
\caption{Two-module physical docking acceptance criteria and outcomes.}
\label{tab:pair-results}
\begin{tabularx}{\textwidth}{L{0.38\textwidth}Y}
\toprule
Criterion & Enforced outcome \\
\midrule
Lifecycle & Exact dock, merge, undock, split semantic sequence \\
Topology & $0\rightarrow1\rightarrow0$ active edges; final two singleton assemblies \\
Failure counters & One dock, one undock, zero failures \\
Capture geometry & Connector translation at latch no greater than approximately \SI{1.001}{mm} \\
Release timing & Undock at least \SI{1.08}{s} after dock (hold plus release delay) \\
Actuation & Both moving wheel positions change by more than \SI{0.2}{rad} \\
Attitude & TILT below \SI{0.03}{rad}; PAN below \SI{0.01}{rad} in the checked final state \\
Separation & Mover reverses to root $x<-0.12$ m while retaining supported height \\
Physics ownership & Root pose/twist/wrench APIs fail the test if used after staging \\
\bottomrule
\end{tabularx}
\end{table}

The experiment establishes that the end-to-end actuation and connection path works for
one deterministic, pre-aligned approach. It does not establish arbitrary-pose docking
robustness.

\subsection{Kinematic Driver-to-Snake demonstration}

The kinematic seven-module demo isolates topology, events, graph behavior, and scenario
orchestration from contact. It stages the initial tree, executes each connector
replacement sequentially, and updates the Runtime Inspector. The topology begins with
seven nodes and six edges, alternates to five and back to six for each action, and ends
in the chain $1-3-2-4-5-6-7$. Gravity and ground are intentionally disabled.

This demonstration remains useful even after the physics scenario exists. It is faster,
deterministic, and can distinguish semantic regressions from control/contact regressions.

\subsection{Physical Driver-to-Snake experiment}

The physical experiment begins with the same six-edge tree staged upright at simulation
time zero. It then disables root motion injection through the regression guard and
executes all four replacements through effort-controlled wheel motion, ground contact,
release delays, measured capture, ideal welds, and paired contact exclusions.

\begin{table}[htbp]
\centering
\caption{Seven-module physical reconfiguration outcomes.}
\label{tab:snake-results}
\begin{tabularx}{\textwidth}{L{0.38\textwidth}Y}
\toprule
Measure & Result or enforced invariant \\
\midrule
Modules & 7 throughout \\
Connection history & $[6,5,6,5,6,5,6,5,6]$ \\
Final topology & Exact chain $1-3-2-4-5-6-7$ with six connector edges \\
Final assemblies & One assembly of size seven \\
Dock outcomes & 10 successful: six initial plus four replacements \\
Undock outcomes & 4 successful \\
Failures & 0 dock and 0 undock failures \\
Assembly events & 10 merges and 4 splits \\
Moving components & $\{1\}$, $\{7\}$, $\{1,2,3\}$, and $\{5,6,7\}$ \\
Component motion & Each wheel in the two three-module movers travels more than \SI{0.5}{rad} \\
Released-face clearance & Greater than \SI{20}{mm} for actions 3 and 4 \\
Completion & Test requires completion by 190 simulated seconds; reference trace about 177 s \\
Final hold & Stable and event-free through 210 simulated seconds \\
Wheel bound & No greater than $\pi/2+0.05\,\mathrm{rad/s}$ in the test \\
Allocation residual & Applied residual no greater than \SI{0.08}{m/s} \\
Pose health & Finite; unit quaternion; up projection $>\cos(0.15)$; root height 0.025--0.06 m \\
Physics ownership & Root pose/twist/wrench APIs forbidden after creation \\
\bottomrule
\end{tabularx}
\end{table}

The semantic event sequence is also exact: six initial
\code{DockCommitted}/\code{AssemblyMerged} pairs followed by four cycles of
\code{UndockCommitted}, \code{AssemblySplit}, \code{DockCommitted}, and
\code{AssemblyMerged}. The final hold verifies that completion does not leave stale
effort, accumulate new events, or destabilize the chain.

\subsection{Runtime visualization experiment}

Runtime Inspector tests use fake clocks and process hosts to check protocol framing,
initialization order, frame delivery, stderr handling, event continuity, graceful stop,
viewer-close behavior, and escalation. Presenter tests verify stable node layout,
parallel-edge curves, selection behavior, and stale-frame rejection. Qt tests run under
an offscreen/Xvfb display.

The current viewer-refresh regression executes ten \SI{2}{ms} physics steps while
synchronizing the native viewer twice, and a second test injects a \SI{30}{ms} sync cost
without causing an immediate repeated refresh. These are scheduling-correctness tests,
not graphics benchmarks.

\subsection{Playback speed}

At factor four, the documented approximately 177-s completion has a nominal wall-clock
target of $177/4\approx44.25$ s, and the 210-s display budget targets 52.5 s. These are
arithmetic targets, not measured achieved rates. The implementation never changes
solver step count: the full seven-module run remains 105,000 nominal \SI{2}{ms} steps.

\subsection{Result interpretation}

The experiments support three conclusions.

\begin{enumerate}
  \item A portable semantic package can drive both mock and physical-runtime workflows
        without moving connection meaning into backend code.
  \item Two-phase commit, canonical events, and derived graph views remain coherent
        through repeated physical dock and release operations.
  \item The backend-neutral joint command path is sufficient to build a nontrivial
        seven-module wheel/contact realization, including connected-component motion.
\end{enumerate}

They do not support claims of route robustness, calibrated hardware accuracy, large-N
scalability, or guaranteed four-times-real-time execution.


\section{Performance, Scalability, and Critical Analysis}
\label{sec:analysis}

\subsection{What has and has not been benchmarked}

The current implementation has strong correctness instrumentation but limited
performance instrumentation. The fresh test run establishes a 537-test, 166.91-second
verification result under coverage on one development environment. It does not isolate
physics step throughput, renderer cost, UI transport latency, memory, or scaling with
module count. Likewise, the documented approximately 177-s simulated completion is a
scenario outcome, not wall-clock performance.

This distinction is important because the demonstration's purpose is presently
functional: prove that the semantic architecture can support wheel/contact dynamics and
reconfiguration. A journal-quality performance claim will require a controlled
benchmark harness, machine disclosure, repeated trials, and raw traces.

\subsection{Algorithmic cost by subsystem}

Let $M$ be module count, $C$ connector count, $E$ active connection count, $J$ joint
count, and $K$ the number of nearby broad-phase pairs.

\begin{table}[htbp]
\centering
\caption{Expected algorithmic costs of principal semantic operations.}
\label{tab:complexity}
\begin{tabularx}{\textwidth}{L{0.28\textwidth}L{0.22\textwidth}Y}
\toprule
Operation & Expected cost & Comment \\
\midrule
Snapshot ingestion & $O(M+J+C)$ & Copies observed link/joint/frame state \\
Spatial hash build/query & expected $O(C+K)$ & Avoids all-pairs $O(C^2)$ in sparse scenes \\
Acceptance/guards per pair & $O(1)$ & Fixed-size vector/quaternion computations \\
Dock commit & component-set dependent & Backend allocation plus assembly merge \\
Undock/split & $O(|V_A|+|E_A|)$ & Recomputes only affected assembly component \\
Topology view build & $O(M+E)$ & Includes isolated nodes and parallel edges \\
Metric summary & $O(M+C+E+|\mathcal L|)$ in simple aggregation & Event history counters can later be incremental \\
Event delta transport & $O(\Delta |\mathcal L|)$ & Sends only unseen canonical events \\
\bottomrule
\end{tabularx}
\end{table}

The broad phase is the primary semantic scaling mechanism. It is effective when
connectors are spatially sparse relative to capture tolerance. Dense piles or very
large tolerance values increase $K$ and approach all-pairs behavior. This has not yet
been measured empirically.

\subsection{Model-view cache behavior}

The model-view factory retains a bounded least-recently-used set of latest snapshots,
defaulting to 32 entries. A topology view token depends on backend sample and topology
revision, so pose motion rebuilds node values while a pure event revision can reuse
graph content. This prevents unbounded per-step graph accumulation.

The cache is correctness tested, but hit rate, build latency, and serialization cost
have not been benchmarked. At seven modules they are negligible relative to physics;
at thousands of modules or multiple simultaneous views they may become visible.

\subsection{Physics and rendering decoupling}

The physical snake default has $210/0.002=105{,}000$ nominal solver steps. A native
viewer synchronization on every step would couple this to 500 attempted visual updates
per simulated second. Commit \code{4a5ca00} caps native synchronization at about 30 Hz,
while retaining every physics step and forcing a final sync. Runtime Inspector graph
publication defaults to 20 Hz wall time, and event deltas remain lossless between
frames.

This is an architectural improvement for accelerated presentation. It reduces
render-call frequency but does not prove a particular achieved real-time factor. The
solver, scene complexity, contacts, Python controller, IPC, native viewer, and host GPU
all contribute. If they take longer than the requested pacing deadline, the runtime
does not skip physics; it simply falls behind the target.

\subsection{Why playback acceleration preserves the experiment}

Changing $\rho$ in \cref{eq:pacing} changes wall waiting only. Contrast three possible
implementations:

\begin{enumerate}
  \item Increasing solver $\Delta t$ reduces step count and changes contact stability.
  \item Multiplying wheel targets changes robot dynamics and event times.
  \item Dividing wall deadlines by $\rho$ preserves simulated trajectories subject to
        deterministic solver execution.
\end{enumerate}

\modsim{} uses the third. Unit tests validate factors 1, 2, 4, and slow motion, reject
zero/non-finite factors, verify protocol round-trip, and inspect viewer deadlines with a
fake clock. A historical protocol record without the new field defaults to 1.0.

\subsection{Physical model debugging lessons}

The physical scenario exposed issues that kinematic staging cannot reveal.

\paragraph{Visual meshes are not contact models.}
Detailed Fusion STLs are suitable for appearance but expensive and fragile for contact.
The pack now uses visual-only STLs and explicit wheel/skid/face primitives.

\paragraph{Rotating square face proxies obstruct locomotion.}
Early LEFT/RIGHT connector proxies were broad square boxes on rotating wheel bodies.
Their corners swept through neighboring module geometry and stalled connected trains.
Coaxial disk/cylinder proxies preserve face contact without creating paddle-like
interference.

\paragraph{Friction direction must follow the wheel.}
Anisotropic coefficients alone are insufficient if the contact tangent stays in a
world-fixed direction. The backend rotates the tangent basis to follow the wheel axle
at every contact solve.

\paragraph{A weld can overconstrain contact.}
Leaving a direct body-pair contact active while also applying a rigid weld causes two
constraints to fight. Pairing each runtime weld with a preallocated exclusion removes
that specific conflict while leaving all other module and assembly contacts active.

\paragraph{Multi-module differential drive is state dependent.}
A controller that moves one module does not automatically move a rigid three-module
train. Assembly allocation, traction, pivot friction, route hysteresis, and lateral
prebias were required. This reinforces the separation between generic semantic plan
and platform-specific physical route.

\subsection{Stability and safety checks}

The physical regressions sample finite poses, quaternions, velocities, and joint state;
bound wheel speed; check upright orientation and supported height; verify released-face
clearance; and hold the final snake beyond completion. Preflight rejects missing effort
support or undeclared limits before staging, and atomic command validation prevents a
partial actuator update.

These are necessary but not sufficient for robust control. They do not measure contact
penetration, energy, tire slip, solver residual, weld impulse, or sensitivity to
parameter perturbation.

\subsection{Determinism and reproducibility}

Semantic ordering is deterministic by construction. Floating-point physics is expected
to be reproducible within one pinned environment and scene, but no trace hash is
currently archived. Exact measured maxima and phase transition timestamps are also not
stored as an artifact. The regression protects thresholds and final state, which is
appropriate for development but insufficient for a comparative experimental paper.

Future result export should record at least:

\begin{itemize}
  \item simulation and wall timestamps;
  \item scenario phase/action/waypoint;
  \item module poses, joint targets/state/effort;
  \item connector position, axis, roll, and speed errors;
  \item contact count, force, penetration, and solver health;
  \item topology, event sequence, and constraint handles; and
  \item software version, dependency lock hash, scene configuration, and random seed.
\end{itemize}

\subsection{Test and CI analysis}

Coverage is branch-aware for the backend-neutral package. Optional \mujoco{} code is
checked separately because generated bindings lack complete stubs; Studio is checked by
its own strict configuration. CI installs only core dependencies in the Python matrix,
\mujoco{} in the headless job, and Studio in the Xvfb job. Package artifacts are built
and smoke-tested only after those jobs.

One current CI gap deserves explicit notice. The three tests in
\pathcode{tests/test_mujoco_contact_mechanics.py} validate authored tire pairs,
wheel-relative friction frames, and atomic weld/exclusion claim/release, but the
dedicated \mujoco{} job's explicit file list omits that file. They ran in the fresh local
537-test suite because \mujoco{} was installed; in a core-only CI job they skip. The file
should be added to the dedicated job before treating contact mechanics as continuously
protected on every pull request.

\subsection{Threats to validity}

\paragraph{Construct validity.}
Graph and event correctness are measured directly against the semantic contract.
Hardware accuracy is not: the contact model and actuators are provisional.

\paragraph{Internal validity.}
The root-motion guard establishes that physical scenarios do not teleport after
creation. Deterministic initial staging and authored routes nevertheless make the task
easier than arbitrary deployment.

\paragraph{External validity.}
Only one physical robot platform and one dynamic backend are exercised. The mock is a
semantic reference, not independent physics validation.

\paragraph{Statistical conclusion validity.}
The physical results are one tuned deterministic regression, not a repeated randomized
experiment. No confidence intervals or success-rate statistics can be claimed.

\paragraph{Artifact validity.}
The source checkout and lock file support reproduction, but there is no committed raw
trace, hardware calibration data set, public license, CITATION metadata, or explicit
public asset license. These are publication blockers for calling the artifact openly
reusable.

\subsection{Recommended benchmark campaign}

A meaningful next performance campaign should include:

\begin{enumerate}
  \item timestep sweep at 1, 2, and 5 ms;
  \item initial gap, lateral error, and yaw perturbation grid;
  \item repeated trace hashes and terminal metrics;
  \item requested versus achieved real-time factor with viewer off/on;
  \item scaling in modules/connectors and broad-phase density;
  \item model-view build/cache/serialization and IPC latency;
  \item contact-exclusion and friction-frame ablations; and
  \item calibrated hardware comparison for wheel and docking trajectories.
\end{enumerate}


\section{Limitations and Future Work}
\label{sec:future}

\subsection{Current limitations}

The present system is an operational vertical slice, not a complete modular-robot
simulation product. The most important limitations are grouped below.

\subsubsection{Physical fidelity}

\begin{itemize}
  \item Tire friction, rear-skid geometry, joint damping/armature, wheel effort limits,
        and feedback gains are provisional.
  \item Detailed visual meshes are not calibrated collision hulls.
  \item EP-Face attraction, magnetic state, electrical switching, and capture forces are
        absent. Accepted connectors transition to an ideal fixed weld.
  \item \mujoco{} equality forces are not reported into core, so runtime connector
        overload and break-force behavior are unverified in physical demos.
  \item Welded chains are solver constraints, not recompiled rigid assemblies, and can
        be softer than an equivalent compiled tree.
  \item The generic \SI{6}{mm} position capture proxy does not reproduce separately
        measured normal and lateral capture envelopes.
\end{itemize}

\subsubsection{Control and planning}

\begin{itemize}
  \item Physical routes are authored and tuned; no autonomous path, obstacle, or
        reconfiguration planner is integrated.
  \item Initial two-module pose is heading aligned, and the seven-module Driver tree is
        staged at time zero.
  \item Only effort commands are implemented by \mujoco{}. Position/velocity commands
        and explicit actuator/transmission models are missing.
  \item Multi-module allocation is a projected rigid-twist heuristic rather than an
        optimal contact-aware controller.
  \item No hardware-in-the-loop or hardware trajectory comparison exists.
\end{itemize}

\subsubsection{Semantic and backend completeness}

\begin{itemize}
  \item Only fixed runtime connections execute in \mujoco{}; compliant, hinge, ball,
        and custom connections remain deferred.
  \item Backend mapping resolution is incomplete, and named-frame-only connectors are
        unsupported in \mujoco{}.
  \item Capability definitions are descriptive rather than executable.
  \item Connector \code{active} is persisted but not yet a runtime enable gate.
  \item Isaac Sim has no adapter. The mock is kinematic and lacks articulated motion.
  \item Physical samples are not event sourced, so the semantic log alone cannot replay
        a continuous trajectory.
\end{itemize}

\subsubsection{Authoring and visualization}

\begin{itemize}
  \item Studio lacks 3-D connector picking/gizmos, joint animation, live assemblies,
        texture support, acceptance overlays, and physics/contact visualization.
  \item The Runtime Inspector renders only one topology graph plus an event table and
        summary status. It lacks pause/restart/scrub, manual control, plots, and a
        general scene editor.
  \item Only one model-view builder is included; plugin discovery and incremental graph
        patches are absent.
  \item Named demos depend on a source-checkout examples tree rather than an installed
        resource package.
\end{itemize}

\subsubsection{Publication and governance}

The repository currently lacks an explicit root software license and CITATION metadata.
The committed mesh metadata records collaborator authorization but is not a public
redistribution license. These issues must be resolved before presenting the source and
SMORES assets as an open, reusable research artifact. Release automation and trusted
package publication are also absent.

\subsection{Phase I: publication-quality validation}

The first future-work phase should improve evidence before adding broad features.

\begin{enumerate}
  \item \textbf{Calibrate the physical plant.} Identify wheel motor/gear dynamics,
        current/torque limits, tire longitudinal/lateral friction, mass distribution,
        inertias, joint friction, and support contact from hardware experiments.
  \item \textbf{Model EP-Face capture.} Add explicit face state and a validated
        magnetic/capture force model before the weld transition. Preserve the ideal
        weld as a selectable fast approximation.
  \item \textbf{Report constraint loads.} Extract \mujoco{} equality forces and map
        them to connector normal, shear, bending, and torsion metrics. Exercise
        overload release against measured thresholds.
  \item \textbf{Archive raw traces.} Export JSON/CSV or a binary time-series format with
        configuration, dependency lock hash, events, state, controller targets, contact,
        and health values.
  \item \textbf{Run sensitivity experiments.} Vary timestep, lateral/yaw error, gap,
        friction, mass, effort, and route perturbations. Report success rates and
        confidence intervals.
  \item \textbf{Compare hardware and simulation.} Measure approach, wheel response,
        capture, release, and reconfiguration trajectories on physical modules and
        quantify error rather than relying on visual similarity.
  \item \textbf{Close verification gaps.} Add the contact-mechanics test file to the
        dedicated \mujoco{} CI job and add a pinned report-build job.
\end{enumerate}

\subsection{Phase II: semantic and backend completeness}

\subsubsection{Backend mappings}

Complete asset, link, joint, connector-frame, and constraint mapping resolution. A
validator should optionally parse selected mechanical assets and cross-check every
mapped source name. Named backend frames should become first-class measured connector
sources.

\subsubsection{Actuator model}

Introduce an actuator/transmission catalog that distinguishes semantic control modes
from motor models. Implement position and velocity modes, rate limits, gear ratio,
torque-speed curves, sensor latency/noise, and multi-rate low-level control. This will
allow the controller period to reflect hardware rather than solver rate.

\subsubsection{Connection types}

Implement compliant constraints by documenting the conversion from physical stiffness
and damping to solver parameters. Add parameterized hinge and ball connections, then a
safe extension mechanism for custom constraints. Conformance tests must verify both
semantic event outcomes and intended degrees of freedom.

\subsubsection{Second dynamic backend}

An Isaac Sim adapter is the natural next cross-backend milestone. It should map pack
assets and connector frames into USD/PhysX, implement runtime fixed constraints and
joint effort, and run the same docking and physical scenario conformance suite. The aim
is not pixel-identical trajectories, but agreement on semantic decisions and stable
identities under each engine's measured state.

\subsection{Phase III: planning and reusable model views}

The current \code{ReconfigurationPlan} already separates topology actions from routes.
Future planners can populate it from graph goals. A complete stack would include:

\begin{itemize}
  \item topology planning and module-to-goal assignment;
  \item collision-aware single-module and connected-assembly motion planning;
  \item helper-module actions for faces that cannot be reached through planar
        differential drive;
  \item closed-loop route replanning and recovery;
  \item concurrent action scheduling with conflict detection; and
  \item capability-aware task planning above reconfiguration.
\end{itemize}

Model views should expand in parallel:

\begin{itemize}
  \item a hardware tree with joints and attached connector ports;
  \item authored and runtime frame graphs;
  \item a bipartite module/port graph;
  \item adjacency, incidence, and ownership matrices;
  \item platform-specific lattice/configuration-coordinate views; and
  \item a hybrid docking view containing candidates, lifecycle, guards, and active
        constraints.
\end{itemize}

All should remain immutable derived DTOs with explicit source revisions. Third-party
discovery can be layered onto the registry without allowing Robot Packs to execute
arbitrary code.

\subsection{Phase IV: authoring and runtime UX}

Studio should gain direct connector glyph selection, a transform gizmo, snapping to CAD
frames/surfaces, visual acceptance volumes, joint configuration/animation, texture
bundling and rendering, and multi-module assembly preview. Pack validation should link
issues to both fields and 3-D entities.

The Runtime Inspector should gain pause/resume, restart, checkpoint and replay, timeline
scrubbing, live connector/acceptance diagnostics, joint controls, time-series metrics,
force/contact panels, and scenario parameter forms. A combined application shell can
host authoring and runtime modes while preserving their separate mutable owners.

\subsection{Phase V: scale, language, and deployment}

Python is appropriate for the present research framework, but larger simulations may
benefit from a C++ stepping/event kernel with Python bindings. The current immutable
DTOs, stable IDs, and strict protocol provide a migration path. Longer-term work may
include:

\begin{itemize}
  \item a C++ runtime for large module counts and deterministic multithreading;
  \item distributed or multi-rate simulation;
  \item ROS 2 and hardware-in-the-loop bridges;
  \item browser clients consuming the versioned immutable frame protocol;
  \item remote compute with local/web visualization; and
  \item a curated registry of independently versioned Robot Packs and scenario data.
\end{itemize}

A browser client need not move physics into the browser. The same separation used by
the current Qt parent and \mujoco{} child can support server-owned simulation and a
remote web frontend.

\subsection{Multi-platform evaluation}

The framework's genericity should be demonstrated, not assumed. At least two additional
platforms should be integrated: one with a different connector gender/constraint model
and one heterogeneous multi-robot system. Useful tests include non-tree topology,
parallel edges, loop closure, non-wheel locomotion, and a connector on an articulated
child link. Only then can claims of platform independence be strengthened beyond
architectural separation.

\subsection{Desired scientific endpoint}

The long-term result should be a calibrated, cross-backend, reproducible experimental
platform where a planner can issue one connector-level plan, different engines can
execute it through their own physics, hardware can consume the same semantic command
stream, and all produce comparable event/topology traces with quantified physical
error. The present work establishes the data and control boundaries needed to pursue
that endpoint.


\section{Conclusion}
\label{sec:conclusion}

Modular-robot simulation requires a coherent account of meaning as well as mechanics.
A physics engine can integrate bodies and solve a weld, but it does not inherently know
which faces are compatible, which orientation was selected, whether policy permits a
dock, when the logical graph should change, or how an interface should explain the
event. \modsim{} supplies that missing semantic layer while deliberately delegating
physics to backend adapters.

The implemented architecture connects a strict portable \robotpack{}, schema-aware
Studio authoring, canonical runtime state, structured docking acceptance, two-phase
physical/logical commit, deterministic events, derived assemblies, atomic joint
commands, backend capability checks, immutable model views, and a dual-window Runtime
Inspector. A dependency-free mock provides executable semantic reference behavior, and
the \mujoco{} adapter adds articulated dynamics, effort actuation, measured connector
sites, contact, runtime welds, paired collision exclusions, and a native viewer.

The \smoresep{} case study demonstrates the value of keeping these layers separate.
The same connector and topology semantics support a fast kinematic lifecycle, a
two-module physical differential-drive dock/undock cycle, and a seven-module
wheel/contact realization of the published Driver-to-Snake connector plan. The physical
reference run performs four $6\rightarrow5\rightarrow6$ topology transitions, reaches
the exact final chain, records ten docks and four undocks without semantic failure, and
does so without root motion injection after initialization. A fresh current-head
verification run passes 537 tests with 86\% branch-aware core coverage.

The strongest conclusion is architectural, not a claim of hardware fidelity: modular
semantics, physical constraints, canonical events, and generated views can remain
coherent while physical fidelity is increased incrementally. The remaining work is
substantial---calibration, magnetic capture, force reporting, autonomous planning,
additional views, richer Studio/runtime interfaces, a second dynamic backend, and
hardware validation---but it can proceed without replacing the foundation established
here.


\appendix
\section{Robot Pack Schema Summary}
\label{app:schema}

This appendix summarizes the implemented format. The normative contract remains
\pathcode{docs/robot_pack_spec.md} and the Pydantic models under
\pathcode{src/modsim/robot_packs/schema.py}.

\begin{longtable}{L{0.24\textwidth}L{0.25\textwidth}L{0.39\textwidth}}
\caption{Robot Pack 0.1 principal documents and fields.}
\label{tab:schema-summary}\\
\toprule
Entity & Principal fields & Purpose \\
\midrule
\endfirsthead
\toprule
Entity & Principal fields & Purpose \\
\midrule
\endhead
Manifest & schema version, ID/name/version, metadata, views, assets, split paths,
mappings & Root identity and cross-document catalog \\
Asset manifest & URDF list, mesh directories, MuJoCo assets, Isaac assets, thumbnails &
Declared mechanical/visual resources \\
Module type & asset reference, root link, mass, joints, connectors, capabilities &
Reusable semantic hardware unit \\
Joint & source name, type, parent/child, axis, control modes, limits &
Semantic articulation and controller contract \\
Connector instance & type, parent link, frame/local pose, docking/approach axes,
metadata & Physical port placement \\
Connector type & activity, gender, compatibility, orientations, acceptance, physical
intent, limits, policy, metadata & Shared mating and lifecycle contract \\
Capability & kind, description, connector requirements & Advertised action/behavior \\
Model-view recipe & builder, modes, default, configuration & Portable view-generation request \\
Backend mapping & asset catalog and module link/joint/frame/constraint maps &
Semantic-to-backend name translation \\
\bottomrule
\end{longtable}

\subsection{Identifier grammar}

Most semantic IDs match
\code{\^{}[a-z][a-z0-9]*(?:\_[a-z0-9]+)*\$} and have a maximum length of 64.
Runtime compound identifiers use \code{<module>/<local>} for connectors/joints and a
canonical sorted pair for connection IDs. Display names are not IDs.

\subsection{Units}

Fields include units in their names wherever ambiguity is likely:
\code{mass\_kg}, \code{xyz\_m}, \code{rpy\_rad},
\code{position\_tolerance\_m}, \code{orientation\_tolerance\_rad},
\code{max\_relative\_velocity\_m\_s}, \code{effort\_nm}, and related forms.
Runtime commands and snapshots use SI.

\subsection{Connector policy defaults}

The effective default is explicit command required, measured alignment, no redock
cooldown, and no break force. Defaults are serialized and round-trip tested so Studio
editing does not silently erase connector-type metadata or policy.

\section{Runtime Events and State Transitions}
\label{app:events}

\begin{longtable}{L{0.27\textwidth}L{0.58\textwidth}}
\caption{Canonical semantic event taxonomy.}
\label{tab:event-taxonomy}\\
\toprule
Event & Meaning \\
\midrule
\endfirsthead
\toprule
Event & Meaning \\
\midrule
\endhead
\code{DockCandidateDetected} & A compatible pair entered candidate evaluation; emitted
on transition rather than every step. \\
\code{DockCommitted} & Backend created a physical constraint and the logical connection
was committed. \\
\code{DockFailed} & An explicit attempt failed compatibility, acceptance, guard, or
backend execution. \\
\code{AssemblyMerged} & A committed dock joined two previously separate components. \\
\code{UndockCommitted} & Backend removed the constraint and logical connection. \\
\code{UndockFailed} & Requested release did not execute. \\
\code{AssemblySplit} & An undock separated one connected component into multiple components. \\
\code{ConnectorOverloaded} & Reported physical load exceeded the weaker break-force policy. \\
\bottomrule
\end{longtable}

The actively used connector lifecycle states include free, candidate detected,
in-acceptance, latching, docked, releasing, and failed. The broader vocabulary also
contains aligning and load-bearing states for future richer execution. A failed
transient is cleared on the next backend ingest, while cooldown prevents immediate
retries where configured.

\subsection{Successful dock pseudocode}

\begin{lstlisting}[language=Python]
snapshot = adapter.snapshot()
world.ingest(snapshot)

for pair in sorted(candidate_pairs(world)):
    compatibility = check_compatibility(pair)
    acceptance = evaluate_acceptance(pair, world)
    guards = evaluate_guards(pair, world, commands)
    if not (compatibility.ok and acceptance.ok and guards.ok):
        report_commanded_failure_if_needed()
        continue

    request = build_connection_request(pair, acceptance)
    outcome = adapter.create_physical_connection(request)
    if outcome.accepted:
        world.apply(DockCommitted(..., handle=outcome.handle))
        world.apply(AssemblyMerged(...))  # only if components differed
    else:
        world.apply(DockFailed(reason="backend_refused", ...))
\end{lstlisting}

\section{SMORES-EP Parameters and Provenance}
\label{app:parameters}

\begin{longtable}{L{0.33\textwidth}L{0.22\textwidth}L{0.31\textwidth}}
\caption{Key \smoresep{} model and controller values.}
\label{tab:parameter-provenance}\\
\toprule
Quantity & Value & Provenance/status \\
\midrule
\endfirsthead
\toprule
Quantity & Value & Provenance/status \\
\midrule
\endhead
Module mass & \SI{0.4386368}{kg} & Fusion-derived geometry/inertial export \\
Wheel radius & \SI{0.040}{m} & Fusion-derived \\
Track width & \SI{0.0672}{m} & Fusion-derived \\
Wheel angular-speed cap & $\pi/2\,\mathrm{rad/s}$ & Published platform value \\
Wheel effort cap & \SI{0.04}{N.m} & Provisional simulation value \\
PAN/TILT effort cap & \SI{0.10}{N.m} & Provisional simulation value \\
Wheel velocity gain & $0.1\,\mathrm{N,m/(rad/s)}$ & Provisional simulation value \\
Hold $P,D$ & $1.0,0.02$ & Provisional simulation value \\
Solver/controller period & \SI{0.002}{s} default & Provisional tuned value \\
Maximum accepted timestep & \SI{0.005}{s} & Scenario safety bound \\
Connector position tolerance & \SI{0.006}{m} & Conservative semantic proxy \\
Pair near-contact gate & \SI{0.001}{m} & Two-module simulation choice \\
Physical-snake latch gate & \SI{0.006}{m} & Uses pack acceptance envelope \\
Orientation tolerance & \SI{25}{\degree} & Literature-motivated, not fully calibrated \\
Relative-speed tolerance & \SI{0.05}{m/s} & Pack semantic value \\
Redock cooldown & \SI{0.1}{s} & Pack semantic value \\
Release delay & \SI{0.08}{s} & Literature-derived switching delay model \\
Normal/shear/bending limits & 100 N / 50 N / 2 N m & Explicitly unverified/provisional \\
Tire friction tuple & $(0.05,1,0.01,10^{-4},10^{-4})$ & Provisional anisotropic contact \\
Skid friction tuple & $(0.08,0.08,0.005,10^{-4},10^{-4})$ & Provisional contact \\
Rigid-allocation residual cap & \SI{0.08}{m/s} & Provisional controller safety bound \\
\bottomrule
\end{longtable}

\section{Driver-to-Snake Topology and Routes}
\label{app:routes}

\subsection{Initial topology}

The six initial connector edges are:

\begin{enumerate}
  \item \code{module\_1/bottom}--\code{module\_2/pan};
  \item \code{module\_2/bottom}--\code{module\_3/pan};
  \item \code{module\_2/right}--\code{module\_4/left};
  \item \code{module\_4/right}--\code{module\_5/left};
  \item \code{module\_5/pan}--\code{module\_6/bottom}; and
  \item \code{module\_5/bottom}--\code{module\_7/pan}.
\end{enumerate}

\subsection{Physical route parameters}

Waypoints are $(f,l,h,d,\epsilon)$: forward offset, left offset, heading offset,
drive direction, and position tolerance, all relative to the action's final moving-root
pose. Heading tolerance is $\pi$ in the current route helper unless explicit alignment
is enabled.

\begin{longtable}{L{0.08\textwidth}L{0.72\textwidth}L{0.10\textwidth}}
\caption{Authored target-relative Driver-to-Snake routes.}
\label{tab:route-details}\\
\toprule
Action & Waypoints $(f,l,h,d,\epsilon)$ & Approach \\
\midrule
\endfirsthead
\toprule
Action & Waypoints $(f,l,h,d,\epsilon)$ & Approach \\
\midrule
\endhead
1 & $(0.34,0,0,+1,0.03)$; $(0.34,0.13,\pi/2,+1,0.03)$;
$(-0.06,0.13,\pi,+1,0.03)$; $(-0.06,0,-\pi/2,+1,0.03)$;
$(-0.03,0,0,+1,0.01)$ & forward \\
2 & $(-0.34,0,0,-1,0.03)$; $(-0.34,-0.13,-\pi/2,+1,0.03)$;
$(0.06,-0.13,0,+1,0.03)$; $(0.06,0,\pi/2,+1,0.03)$;
$(0.12,0,0,+1,0.015)$ & reverse \\
3 & $(-0.27,0.088,0,-1,0.03)$; $(-0.03,0,-0.35,+1,0.01)$ & forward \\
4 & $(0.27,-0.088,0,+1,0.075)$; $(0.03,0.008,-0.35,-1,0.01)$ & reverse \\
\bottomrule
\end{longtable}

The action-4 \SI{8}{mm} prebias and \SI{75}{mm} clearance tolerance are empirical
simulation tuning. They should be replaced by planned, feedback-corrected paths in a
general controller.

\section{Reproduction Commands}
\label{app:commands}

\subsection{Environment}

\begin{lstlisting}[language=bash]
uv sync --locked --extra dev --extra studio --extra mujoco --python 3.12
uv run --no-sync modsim --version
\end{lstlisting}

Equivalent editable pip installation:

\begin{lstlisting}[language=bash]
python3.12 -m venv .venv
.venv/bin/python -m pip install -e ".[studio,mujoco,dev]"
\end{lstlisting}

\subsection{Pack inspection and Studio}

\begin{lstlisting}[language=bash]
modsim pack inspect examples/robot_packs/smores_ep
modsim pack validate examples/robot_packs/smores_ep
modsim pack validate examples/robot_packs/smores_ep --profile simulation
modsim studio examples/robot_packs/smores_ep
\end{lstlisting}

\subsection{Physical two-module demonstration}

\begin{lstlisting}[language=bash]
modsim runtime examples/robot_packs/smores_ep \
  --backend mujoco \
  --demo smores_diff_drive_dock_undock \
  --model-view smores_topology
\end{lstlisting}

\subsection{Kinematic Driver-to-Snake}

\begin{lstlisting}[language=bash]
modsim runtime examples/robot_packs/smores_ep \
  --backend mujoco \
  --demo smores_driver_to_snake \
  --model-view smores_topology \
  --duration 14 \
  --dt 0.002 \
  --no-gravity
\end{lstlisting}

\subsection{Physical Driver-to-Snake}

\begin{lstlisting}[language=bash]
modsim runtime examples/robot_packs/smores_ep \
  --backend mujoco \
  --demo smores_physical_driver_to_snake \
  --model-view smores_topology \
  --speed 4
\end{lstlisting}

Omit \code{--speed} for real-time pacing. The factor is fixed at launch. Both windows
retain the final state until Stop or close.

\subsection{Verification}

\begin{lstlisting}[language=bash]
uv lock --check
uv run --no-sync ruff check .
uv run --no-sync ruff format --check .
uv run --no-sync pyright
uv run --no-sync pyright --project pyright-mujoco.json
uv run --no-sync pyright --project pyright-studio.json
uv run --no-sync pytest
\end{lstlisting}

Focused physical tests:

\begin{lstlisting}[language=bash]
uv run --no-sync pytest tests/test_mujoco_physical_docking.py
uv run --no-sync pytest tests/test_mujoco_physical_reconfiguration.py
uv run --no-sync pytest tests/test_mujoco_contact_mechanics.py
\end{lstlisting}

\section{Requirement-to-Evidence Traceability}
\label{app:traceability}

\begin{longtable}{L{0.11\textwidth}L{0.35\textwidth}L{0.40\textwidth}}
\caption{Traceability from \cref{tab:requirements} to principal implementation evidence.}
\label{tab:traceability}\\
\toprule
Req. & Implementation & Principal regression areas \\
\midrule
\endfirsthead
\toprule
Req. & Implementation & Principal regression areas \\
\midrule
\endhead
R1--R2 & schema, loader, validator & schema, loader, validator, SMORES example pack \\
R3 & backend base/registry, mock, MuJoCo adapter & registry, backend conformance \\
R4 & broad phase, acceptance, guards, docking & acceptance and docking suites \\
R5 & connection request/outcome and event commit & docking, MuJoCo weld, refusal tests \\
R6 & world revisions, assemblies, model topology & assemblies, revisions, model views \\
R7 & joint IDs/commands/session validation, MuJoCo effort & joint commands, MuJoCo backend \\
R8 & inspection DTO/protocol/presenter/process/window & runtime inspection/process/widgets/worker \\
R9 & safe paths, writer staging/rollback, Studio project & loader, writer, Studio project \\
R10 & backend and model-view registries & registry and model-view factory tests \\
Physical pair & differential-drive scenario and MJCF contact & physical docking/contact tests \\
Physical snake & physical reconfiguration and authored routes & full MuJoCo physical reconfiguration \\
\bottomrule
\end{longtable}

\section{Artifact Checklist}
\label{app:artifact}

For an archival release accompanying a future publication, include:

\begin{itemize}
  \item tagged commit, version, lock file, and built wheel/source distribution;
  \item explicit software LICENSE, asset license/NOTICE, and CITATION.cff;
  \item pinned TeX toolchain and rendered report PDF;
  \item exact command configurations and source-checkout scenarios;
  \item raw trace files with hashes;
  \item machine/OS/CPU/GPU disclosure in a separate benchmark manifest;
  \item videos synchronized with event sequence and simulated time;
  \item calibration data and hardware experimental protocol; and
  \item expected terminal metrics and tolerance bands.
\end{itemize}



\bibliographystyle{plainnat}
\bibliography{references}

\end{document}
